% =========================================================
% Chapter: Current, Resistance, Capacitance, and DC Circuits
% POSN 1 Physics Preparation Notes
% Language: Thai
% Compiler: XeLaTeX
% This file assumes the main project already defines:
%   examplebox, exercisebox, tcolorbox, siunitx, graphicx
% =========================================================

\chapter{วงจรไฟฟ้ากระแสตรง ตัวต้านทาน และตัวเก็บประจุ}

\begin{center}
    {\Large \textbf{Current, Resistance, Capacitance, and Direct-Current Circuits}}\\[4pt]
    {\normalsize เอกสารเตรียมสอบคัดเลือก สอวน. ฟิสิกส์ ค่าย 1}\\[6pt]
    {\normalsize จัดทำโดย \textbf{Tames Thanakrit Damduan}}\\[2pt]
    {\footnotesize \textcopyright\ 2026 Tames Thanakrit Damduan. All rights reserved.}\\[8pt]
\end{center}

\section*{เป้าหมายของบทเรียน}

หลังจากเรียนบทนี้ นักเรียนควรทำสิ่งต่อไปนี้ได้

\begin{enumerate}
    \item อธิบายกระแสไฟฟ้าจากอัตราการไหลของประจุ และใช้ $I=dq/dt$ ได้
    \item เชื่อมโยงกฎของโอห์มกับสภาพต้านทานของวัสดุ และอนุมาน $R=\rho L/A$ ได้
    \item อนุมานกำลังไฟฟ้า $P=I\Delta V$, $P=I^2R$ และ $P=(\Delta V)^2/R$ ได้
    \item วิเคราะห์แหล่งกำเนิดแรงเคลื่อนไฟฟ้าที่มีความต้านทานภายในได้
    \item อนุมานสูตรตัวต้านทานต่ออนุกรมและขนานจากการอนุรักษ์ประจุและพลังงานได้
    \item ใช้กฎของเคอร์ชอฟฟ์วิเคราะห์วงจรหลายแขนงและหลายลูปได้
    \item อธิบายความจุไฟฟ้า การต่อตัวเก็บประจุ และพลังงานในตัวเก็บประจุได้
    \item วิเคราะห์สภาวะเริ่มต้น สภาวะคงตัว และการเปลี่ยนแปลงตามเวลาในวงจร $RC$ ได้
    \item ใช้สมมาตร จุดศักย์เท่ากัน และการลดวงจรแบบย้อนกลับแก้โจทย์วงจรซับซ้อนได้
\end{enumerate}

\begin{tcolorbox}[title={แนวคิดสำคัญของบทนี้}]
สมการของวงจรไม่ได้เป็นชุดสูตรที่ต้องจำแยกจากกัน แต่เกิดจากหลักพื้นฐานเพียงไม่กี่ข้อ

\begin{enumerate}
    \item การอนุรักษ์ประจุ ให้กฎกระแสที่จุดต่อ
    \item การอนุรักษ์พลังงาน ให้กฎความต่างศักย์รอบลูป
    \item ความสัมพันธ์ของอุปกรณ์ เช่น $\Delta V=IR$ และ $q=C\Delta V$
\end{enumerate}

เมื่อระบุจุดที่มีศักย์เท่ากันและเขียนสามหลักนี้อย่างเป็นระบบ วงจรที่ดูซับซ้อนจะกลายเป็นระบบสมการธรรมดา
\end{tcolorbox}

% =========================================================
\section{กระแสไฟฟ้า}

\subsection{นิยามของกระแสไฟฟ้า}

ให้ประจุสุทธิ $\Delta q$ ผ่านพื้นที่หน้าตัดของตัวนำในช่วงเวลา $\Delta t$ นิยามกระแสเฉลี่ยเป็น

\[
I_{\text{avg}}\equiv\frac{\Delta q}{\Delta t}
\]

ถ้าต้องการกระแส ณ ขณะหนึ่ง ให้ลดช่วงเวลาเข้าใกล้ศูนย์

\[
I\equiv\lim_{\Delta t\to 0}\frac{\Delta q}{\Delta t}
\]

ดังนั้น

\begin{tcolorbox}[title={นิยามกระแสไฟฟ้า}]
\[
\boxed{I=\frac{dq}{dt}}
\]

หน่วย SI คือแอมแปร์

\[
1\,\si{A}=1\,\si{C.s^{-1}}
\]
\end{tcolorbox}

ทิศกระแสตามข้อตกลงคือทิศที่ประจุบวกเคลื่อนที่ ในโลหะตัวพาหะประจุส่วนใหญ่เป็นอิเล็กตรอน จึงเคลื่อนที่สวนทางกับทิศกระแสตามข้อตกลง

\subsection{กระแสคงตัวและปริมาณประจุที่ผ่าน}

จากนิยาม

\[
I=\frac{dq}{dt}
\]

จัดรูปได้

\[
dq=I\,dt
\]

อินทิเกรตตั้งแต่ $t_1$ ถึง $t_2$

\[
\Delta q=\int_{t_1}^{t_2}I(t)\,dt
\]

ถ้ากระแสคงตัว

\[
\boxed{\Delta q=I\Delta t}
\]

ถ้าตัวพาประจุแต่ละตัวมีขนาดประจุ $e$ จำนวนอนุภาคที่ผ่านหน้าตัดคือ

\[
N=\frac{\Delta q}{e}=\frac{I\Delta t}{e}
\]

\subsection{แบบจำลองการนำไฟฟ้าและความเร็วลอยเลื่อน}

พิจารณาตัวนำพื้นที่หน้าตัด $A$ มีตัวพาประจุจำนวน $n$ ตัวต่อหนึ่งหน่วยปริมาตร แต่ละตัวมีประจุขนาด $q_c$ และมีความเร็วลอยเลื่อนขนาด $v_d$

ในเวลา $\Delta t$ ตัวพาประจุที่ผ่านหน้าตัดมาจากทรงกระบอกยาว

\[
\Delta x=v_d\Delta t
\]

ปริมาตรของทรงกระบอกคือ

\[
\Delta V=A v_d\Delta t
\]

จำนวนตัวพาประจุในปริมาตรนี้คือ

\[
N=nAv_d\Delta t
\]

ประจุรวมที่ผ่านหน้าตัดจึงเป็น

\[
\Delta q=Nq_c=nq_cAv_d\Delta t
\]

หารด้วย $\Delta t$

\begin{tcolorbox}[title={กระแสจากแบบจำลองจุลภาค}]
\[
\boxed{I=nq_cAv_d}
\]
\end{tcolorbox}

ความหนาแน่นกระแสมีนิยามเป็นกระแสต่อพื้นที่หน้าตัด

\[
J\equiv\frac{I}{A}
\]

จึงได้

\[
\boxed{J=nq_cv_d}
\]

% =========================================================
\section{ความต้านทานและสภาพต้านทานไฟฟ้า}

\subsection{ความสัมพันธ์แบบโอห์ม}

สำหรับตัวนำโอห์มมิกที่อุณหภูมิคงตัว การทดลองพบว่าความต่างศักย์แปรผันตรงกับกระแส

\[
\Delta V\propto I
\]

ค่าคงตัวของการแปรผันเรียกว่า ความต้านทาน

\begin{tcolorbox}[title={นิยามความต้านทานของอุปกรณ์โอห์มมิก}]
\[
\boxed{R\equiv\frac{\Delta V}{I}}
\]

หรือ

\[
\boxed{\Delta V=IR}
\]

หน่วยคือโอห์ม

\[
1\,\si{\ohm}=1\,\si{V.A^{-1}}
\]
\end{tcolorbox}

สมการ $\Delta V=IR$ ไม่ใช่กฎสากลของอุปกรณ์ทุกชนิด อุปกรณ์ที่กราฟ $I$ กับ $\Delta V$ ไม่เป็นเส้นตรง เช่น ไดโอด ไม่สามารถใช้ค่า $R$ คงตัวเพียงค่าเดียวได้

\subsection{อนุมานความต้านทานของลวดสม่ำเสมอ}

สำหรับวัสดุโอห์มมิก ความสัมพันธ์ระดับจุลภาคระหว่างสนามไฟฟ้าและความหนาแน่นกระแสเขียนเป็น

\[
\vec{E}=\rho\vec{J}
\]

โดย $\rho$ คือสภาพต้านทานไฟฟ้าของวัสดุ

สำหรับลวดยาว $L$ พื้นที่หน้าตัด $A$ และสนามสม่ำเสมอ

\[
\Delta V=EL
\]

และ

\[
J=\frac{I}{A}
\]

แทน $E=\rho J$

\[
\Delta V=(\rho J)L
\]

\[
\Delta V=\rho\left(\frac{I}{A}\right)L
\]

\[
\Delta V=I\left(\rho\frac{L}{A}\right)
\]

เปรียบเทียบกับ $\Delta V=IR$

\begin{tcolorbox}[title={ความต้านทานของลวดสม่ำเสมอ}]
\[
\boxed{R=\rho\frac{L}{A}}
\]
\end{tcolorbox}

สมการนี้แสดงว่า

\[
R\propto L
\]

และ

\[
R\propto\frac{1}{A}
\]

หน่วยของ $\rho$ คือ

\[
\si{\ohm.m}
\]

สภาพนำไฟฟ้า $\sigma$ นิยามเป็นส่วนกลับของสภาพต้านทาน

\[
\boxed{\sigma\equiv\frac{1}{\rho}}
\]

จึงเขียนได้ว่า

\[
\vec{J}=\sigma\vec{E}
\]

\subsection{การเปลี่ยนรูปร่างของลวดเมื่อปริมาตรคงตัว}

ถ้าลวดเดิมยาว $L$ พื้นที่หน้าตัด $A$ และถูกดึงให้ยาวเป็น $xL$ โดยปริมาตรคงตัว

\[
AL=A'(xL)
\]

ดังนั้น

\[
A'=\frac{A}{x}
\]

ความต้านทานใหม่คือ

\[
R'=\rho\frac{xL}{A/x}
\]

\[
\boxed{R'=x^2R}
\]

ผลนี้สำคัญมาก เพราะการดึงลวดไม่ได้เพิ่มเพียงความยาว แต่ยังลดพื้นที่หน้าตัดพร้อมกัน

\subsection{ความต้านทานกับอุณหภูมิ}

ในช่วงอุณหภูมิไม่กว้างมาก สภาพต้านทานของโลหะประมาณได้ด้วยความสัมพันธ์เชิงเส้น

\[
\rho=\rho_0\left[1+\alpha(T-T_0)\right]
\]

ถ้ารูปร่างของลวดเปลี่ยนน้อยมากจนถือว่า $L/A$ คงตัว จาก

\[
R=\rho\frac{L}{A}
\]

จึงได้

\begin{tcolorbox}[title={ความต้านทานกับอุณหภูมิ}]
\[
\boxed{R=R_0\left[1+\alpha(T-T_0)\right]}
\]
\end{tcolorbox}

สมการนี้เป็นแบบจำลองประมาณ ไม่ควรนำไปใช้ไกลจากช่วงอุณหภูมิที่กำหนด

\begin{examplebox}{ตัวอย่าง: ลวดถูกดึงโดยปริมาตรคงตัว}
ลวดเส้นหนึ่งมีความต้านทานเดิม $R$ ถูกดึงให้ยาวเป็น $3$ เท่าของเดิม โดยวัสดุและปริมาตรไม่เปลี่ยน จงหาความต้านทานใหม่

เพราะ

\[
R'=x^2R
\]

เมื่อ $x=3$

\[
\boxed{R'=9R}
\]
\end{examplebox}

% =========================================================
\section{พลังงานและกำลังไฟฟ้า}

\subsection{อนุมานกำลังไฟฟ้าจากพลังงานศักย์ไฟฟ้า}

เมื่อประจุ $dq$ เคลื่อนผ่านความต่างศักย์ $\Delta V$ การเปลี่ยนพลังงานมีขนาด

\[
dU=\Delta V\,dq
\]

นิยามกำลังคืออัตราการถ่ายโอนพลังงาน

\[
P\equiv\frac{dU}{dt}
\]

แทน $dU=\Delta V\,dq$

\[
P=\Delta V\frac{dq}{dt}
\]

และจากนิยามกระแส

\[
I=\frac{dq}{dt}
\]

จึงได้

\begin{tcolorbox}[title={กำลังไฟฟ้า}]
\[
\boxed{P=I\Delta V}
\]
\end{tcolorbox}

หน่วยคือวัตต์

\[
1\,\si{W}=1\,\si{J.s^{-1}}
\]

\subsection{กำลังที่เปลี่ยนเป็นพลังงานภายในของตัวต้านทาน}

สำหรับตัวต้านทานโอห์มมิก

\[
\Delta V=IR
\]

แทนใน $P=I\Delta V$

\[
P=I(IR)
\]

จึงได้

\[
\boxed{P=I^2R}
\]

อีกทางหนึ่ง จาก

\[
I=\frac{\Delta V}{R}
\]

แทนใน $P=I\Delta V$

\[
P=\left(\frac{\Delta V}{R}\right)\Delta V
\]

จึงได้

\begin{tcolorbox}[title={กำลังที่ตัวต้านทานรับ}]
\[
\boxed{P=I\Delta V=I^2R=\frac{(\Delta V)^2}{R}}
\]
\end{tcolorbox}

ต้องเลือกใช้รูปให้สอดคล้องกับสิ่งที่คงตัว

\begin{itemize}
    \item ถ้ากระแสคงตัว ใช้ $P=I^2R$ จะเห็นว่า $P$ เพิ่มตาม $R$
    \item ถ้าความต่างศักย์คงตัว ใช้ $P=(\Delta V)^2/R$ จะเห็นว่า $P$ ลดเมื่อ $R$ เพิ่ม
\end{itemize}

พลังงานที่เปลี่ยนเป็นพลังงานภายในในเวลา $t$ เมื่อกำลังคงตัวคือ

\[
\boxed{E=Pt}
\]

\subsection{พลังงานจูลในตัวต้านทาน}

เมื่อกระแสคงตัว

\[
E_R=P_Rt
\]

ดังนั้น

\[
\boxed{E_R=I^2Rt=I\Delta Vt=\frac{(\Delta V)^2}{R}t}
\]

คำว่า ``กำลังสูญเสีย'' หมายถึงพลังงานไฟฟ้าถูกเปลี่ยนเป็นพลังงานภายในหรือความร้อน ไม่ได้หมายถึงพลังงานหายไปจากกฎอนุรักษ์พลังงาน

\begin{examplebox}{ตัวอย่าง: กำลังหลังดึงลวด}
ลวดต่อกับแหล่งจ่ายที่รักษาความต่างศักย์คงตัว ถูกดึงให้ยาวเป็น $x$ เท่าของเดิมโดยปริมาตรคงตัว จงหาอัตราส่วนกำลังใหม่ต่อกำลังเดิม

จากหัวข้อก่อนหน้า

\[
R'=x^2R
\]

เมื่อความต่างศักย์คงตัว

\[
P=\frac{V^2}{R}
\]

ดังนั้น

\[
\frac{P'}{P}
=\frac{V^2/(x^2R)}{V^2/R}
\]

\[
\boxed{\frac{P'}{P}=\frac{1}{x^2}}
\]
\end{examplebox}

% =========================================================
\section{แรงเคลื่อนไฟฟ้าและความต้านทานภายใน}

\subsection{นิยามแรงเคลื่อนไฟฟ้า}

แรงเคลื่อนไฟฟ้าไม่ใช่แรง แต่เป็นพลังงานที่แหล่งกำเนิดจ่ายต่อหนึ่งหน่วยประจุ

\begin{tcolorbox}[title={นิยามแรงเคลื่อนไฟฟ้า}]
\[
\boxed{\mathcal{E}\equiv\frac{dW_{\text{source}}}{dq}}
\]
\end{tcolorbox}

หน่วยของ $\mathcal{E}$ คือโวลต์

\subsection{แรงดันคร่อมขั้วของแบตเตอรี่จริง}

ให้แบตเตอรี่มีแรงเคลื่อนไฟฟ้า $\mathcal{E}$ และความต้านทานภายใน $r$ ขณะจ่ายกระแส $I$

ใช้การอนุรักษ์พลังงานต่อหนึ่งหน่วยประจุ

\[
\mathcal{E}=V_{\text{terminal}}+Ir
\]

ดังนั้น

\begin{tcolorbox}[title={แบตเตอรี่กำลังจ่ายกระแส}]
\[
\boxed{V_{\text{terminal}}=\mathcal{E}-Ir}
\]
\end{tcolorbox}

ถ้ากระแสไหลเข้าสู่ขั้วบวกเพื่อชาร์จแบตเตอรี่

\[
\boxed{V_{\text{terminal}}=\mathcal{E}+Ir}
\]

ถ้าต่อโหลดภายนอก $R$ เพียงตัวเดียว

\[
\mathcal{E}=I(R+r)
\]

จึงได้

\[
\boxed{I=\frac{\mathcal{E}}{R+r}}
\]

\subsection{กำลังสูงสุดที่ส่งให้โหลด}

กำลังที่โหลด $R$ รับคือ

\[
P_R=I^2R
\]

แทน

\[
I=\frac{\mathcal{E}}{R+r}
\]

ได้

\[
P_R=\frac{\mathcal{E}^2R}{(R+r)^2}
\]

ใช้ความจริงว่า

\[
(R+r)^2\geq 4Rr
\]

จึงมี

\[
P_R\leq\frac{\mathcal{E}^2R}{4Rr}
\]

\[
P_R\leq\frac{\mathcal{E}^2}{4r}
\]

เครื่องหมายเท่ากับเกิดเมื่อ

\[
R=r
\]

ดังนั้น

\begin{tcolorbox}[title={เงื่อนไขส่งกำลังสูงสุดให้โหลด}]
\[
\boxed{R=r}
\]

\[
\boxed{P_{\max}=\frac{\mathcal{E}^2}{4r}}
\]
\end{tcolorbox}

% =========================================================
\section{การต่อตัวต้านทาน}

\subsection{ตัวต้านทานต่ออนุกรม}

ตัวต้านทานต่ออนุกรมมีเส้นทางกระแสเพียงเส้นทางเดียว จึงมีกระแสเท่ากันทุกตัว

\[
I_1=I_2=\cdots=I
\]

ความต่างศักย์รวมเท่ากับผลรวมของความต่างศักย์แต่ละตัว

\[
\Delta V=\Delta V_1+\Delta V_2+\cdots
\]

ใช้ $\Delta V_i=IR_i$

\[
\Delta V=IR_1+IR_2+\cdots
\]

\[
\Delta V=I(R_1+R_2+\cdots)
\]

แต่ความต้านทานสมมูลนิยามจาก

\[
\Delta V=IR_{\text{eq}}
\]

จึงได้

\begin{tcolorbox}[title={ตัวต้านทานต่ออนุกรม}]
\[
\boxed{R_{\text{eq}}=R_1+R_2+R_3+\cdots}
\]
\end{tcolorbox}

ความต้านทานสมมูลแบบอนุกรมต้องมากกว่าความต้านทานทุกตัวในชุด

\subsection{ตัวต้านทานต่อขนาน}

ตัวต้านทานต่อขนานต่ออยู่ระหว่างจุดสองจุดเดียวกัน จึงมีความต่างศักย์เท่ากัน

\[
\Delta V_1=\Delta V_2=\cdots=\Delta V
\]

จากการอนุรักษ์ประจุที่จุดต่อ

\[
I=I_1+I_2+\cdots
\]

ใช้

\[
I_i=\frac{\Delta V}{R_i}
\]

จึงได้

\[
I=\frac{\Delta V}{R_1}+\frac{\Delta V}{R_2}+\cdots
\]

\[
I=\Delta V\left(\frac{1}{R_1}+\frac{1}{R_2}+\cdots\right)
\]

แต่

\[
I=\frac{\Delta V}{R_{\text{eq}}}
\]

ดังนั้น

\begin{tcolorbox}[title={ตัวต้านทานต่อขนาน}]
\[
\boxed{\frac{1}{R_{\text{eq}}}
=\frac{1}{R_1}+\frac{1}{R_2}+\frac{1}{R_3}+\cdots}
\]
\end{tcolorbox}

สำหรับสองตัว

\[
\boxed{R_{\text{eq}}=\frac{R_1R_2}{R_1+R_2}}
\]

ความต้านทานสมมูลแบบขนานต้องน้อยกว่าความต้านทานที่น้อยที่สุดในชุด

\subsection{ตัวแบ่งแรงดัน}

ถ้า $R_1$ และ $R_2$ ต่ออนุกรมกับแหล่งจ่าย $V$

\[
I=\frac{V}{R_1+R_2}
\]

แรงดันคร่อม $R_2$ คือ

\[
V_2=IR_2
\]

แทน $I$

\begin{tcolorbox}[title={ตัวแบ่งแรงดัน}]
\[
\boxed{V_2=V\frac{R_2}{R_1+R_2}}
\]
\end{tcolorbox}

\subsection{ตัวแบ่งกระแสสำหรับตัวต้านทานสองตัว}

ถ้า $R_1$ และ $R_2$ ต่อขนานและมีกระแสรวม $I$ เข้าจุดต่อ ทั้งสองแขนงมีแรงดันเท่ากัน

\[
I_1R_1=I_2R_2
\]

และ

\[
I=I_1+I_2
\]

แก้สมการได้

\[
\boxed{I_1=I\frac{R_2}{R_1+R_2}}
\]

\[
\boxed{I_2=I\frac{R_1}{R_1+R_2}}
\]

กระแสจะแบ่งไปทางแขนงที่มีความต้านทานน้อยกว่าในสัดส่วนที่มากกว่า

\begin{examplebox}{ตัวอย่าง: ตัวต้านทานสามตัว}
ตัวต้านทาน $R$, $R$ และ $2R$ ต้องใช้ครบทุกตัว จงหาอัตราส่วนระหว่างความต้านทานรวมที่มากที่สุดกับน้อยที่สุด

ค่ามากที่สุดเกิดจากต่ออนุกรมทั้งหมด

\[
R_{\max}=R+R+2R=4R
\]

ค่าน้อยที่สุดเกิดจากต่อขนานทั้งหมด

\[
\frac{1}{R_{\min}}
=\frac{1}{R}+\frac{1}{R}+\frac{1}{2R}
=\frac{5}{2R}
\]

\[
R_{\min}=\frac{2R}{5}
\]

ดังนั้น

\[
\boxed{\frac{R_{\max}}{R_{\min}}=10}
\]
\end{examplebox}

% =========================================================
\section{กฎของเคอร์ชอฟฟ์}

\subsection{กฎกระแสที่จุดต่อ}

ถ้าไม่มีประจุสะสมที่จุดต่อ อัตราที่ประจุไหลเข้าต้องเท่ากับอัตราที่ประจุไหลออก

\[
\sum I_{\text{in}}=\sum I_{\text{out}}
\]

หรือเมื่อกำหนดเครื่องหมายให้กระแสเข้าเป็นบวกและออกเป็นลบ

\begin{tcolorbox}[title={กฎจุดต่อของเคอร์ชอฟฟ์}]
\[
\boxed{\sum_{\text{junction}} I=0}
\]
\end{tcolorbox}

กฎนี้เป็นผลโดยตรงจากการอนุรักษ์ประจุ

\subsection{กฎความต่างศักย์รอบลูป}

เมื่อประจุเคลื่อนครบหนึ่งรอบและกลับมาที่จุดเดิม พลังงานศักย์ต่อหนึ่งหน่วยประจุต้องกลับมาเท่าเดิม

\begin{tcolorbox}[title={กฎลูปของเคอร์ชอฟฟ์}]
\[
\boxed{\sum_{\text{closed loop}}\Delta V=0}
\]
\end{tcolorbox}

กฎนี้เป็นผลจากการอนุรักษ์พลังงาน

\subsection{กฎเครื่องหมายที่ควรใช้ให้สม่ำเสมอ}

เมื่อเดินรอบลูป

\begin{enumerate}
    \item ผ่านตัวต้านทานตามทิศกระแส ใช้ $-IR$
    \item ผ่านตัวต้านทานสวนทิศกระแส ใช้ $+IR$
    \item ผ่านแหล่งกำเนิดจากขั้วลบไปขั้วบวก ใช้ $+\mathcal{E}$
    \item ผ่านแหล่งกำเนิดจากขั้วบวกไปขั้วลบ ใช้ $-\mathcal{E}$
\end{enumerate}

ถ้าคำนวณกระแสได้ค่าลบ แปลว่าทิศจริงตรงข้ามกับทิศที่สมมติ ไม่ได้แปลว่าคำตอบผิด

\subsection{ขั้นตอนแก้วงจรหลายลูป}

\begin{enumerate}
    \item กำหนดทิศกระแสในแต่ละแขนงโดยเลือกได้ตามสะดวก
    \item เขียนกฎจุดต่อสำหรับจุดต่ออิสระ
    \item เลือกลูปอิสระและเขียนกฎลูป
    \item แก้ระบบสมการพร้อมกัน
    \item ตรวจหน่วย เครื่องหมาย และกฎอนุรักษ์กำลัง
\end{enumerate}

\subsection{การตรวจคำตอบด้วยกำลัง}

ในวงจรสภาวะคงตัว ผลรวมกำลังที่แหล่งกำเนิดจ่ายควรเท่ากับผลรวมกำลังที่อุปกรณ์รับ

\[
\sum P_{\text{source}}=\sum P_{\text{absorbed}}
\]

สำหรับตัวต้านทานทุกตัว

\[
P_R=I^2R\geq 0
\]

% =========================================================
\section{เครื่องมือวัดในวงจร}

\subsection{แอมมิเตอร์}

แอมมิเตอร์ใช้วัดกระแส จึงต้องต่ออนุกรมกับแขนงที่ต้องการวัด

แอมมิเตอร์อุดมคติมีความต้านทานเป็นศูนย์

\[
R_A=0
\]

เพื่อไม่ให้เปลี่ยนกระแสเดิมของวงจร

\subsection{โวลต์มิเตอร์}

โวลต์มิเตอร์ใช้วัดความต่างศักย์ จึงต้องต่อขนานระหว่างจุดสองจุด

โวลต์มิเตอร์อุดมคติมีความต้านทานเป็นอนันต์

\[
R_V\to\infty
\]

เพื่อไม่ดึงกระแสจากวงจร

\subsection{โอห์มมิเตอร์}

การวัดความต้านทานสมมูลด้วยโอห์มมิเตอร์ควรตัดแหล่งจ่ายภายนอกออกก่อน เพราะโอห์มมิเตอร์มีแหล่งจ่ายภายในของตัวเอง และค่าที่อ่านได้ขึ้นกับเส้นทางทั้งหมดที่เชื่อมระหว่างขั้ววัด

% =========================================================
\section{ตัวเก็บประจุและความจุไฟฟ้า}

\subsection{นิยามความจุไฟฟ้า}

ตัวเก็บประจุประกอบด้วยตัวนำสองส่วนที่มีประจุขนาดเท่ากันแต่เครื่องหมายตรงข้าม

\[
+q\qquad\text{และ}\qquad -q
\]

นิยามความจุไฟฟ้าเป็น

\begin{tcolorbox}[title={นิยามความจุไฟฟ้า}]
\[
\boxed{C\equiv\frac{q}{\Delta V}}
\]

หรือ

\[
\boxed{q=C\Delta V}
\]
\end{tcolorbox}

หน่วยคือฟารัด

\[
1\,\si{F}=1\,\si{C.V^{-1}}
\]

สำหรับตัวเก็บประจุเชิงเส้น ค่า $C$ ขึ้นกับรูปทรง ขนาด ระยะห่าง และสสารระหว่างตัวนำ แต่ไม่ขึ้นกับ $q$ หรือ $\Delta V$ โดยแยกกัน

\subsection{อนุมานความจุของตัวเก็บประจุแผ่นขนาน}

ตัวเก็บประจุแผ่นขนานมีพื้นที่แผ่น $A$ ระยะห่าง $d$ และละผลขอบ

ความหนาแน่นประจุผิวคือ

\[
\sigma=\frac{q}{A}
\]

จากกฎของเกาส์ สนามระหว่างแผ่นคือ

\[
E=\frac{\sigma}{\varepsilon_0}
\]

แทน $\sigma=q/A$

\[
E=\frac{q}{\varepsilon_0A}
\]

สนามสม่ำเสมอให้ความต่างศักย์

\[
\Delta V=Ed
\]

ดังนั้น

\[
\Delta V=\frac{q}{\varepsilon_0A}d
\]

จากนิยาม $C=q/\Delta V$

\[
C=\frac{q}{qd/(\varepsilon_0A)}
\]

จึงได้

\begin{tcolorbox}[title={ตัวเก็บประจุแผ่นขนาน}]
\[
\boxed{C=\varepsilon_0\frac{A}{d}}
\]
\end{tcolorbox}

ถ้าเติมไดอิเล็กทริกเต็มช่องว่างและมีค่าคงตัวไดอิเล็กทริก $\kappa$

\[
\varepsilon=\kappa\varepsilon_0
\]

จึงได้

\[
\boxed{C=\kappa\varepsilon_0\frac{A}{d}=\kappa C_0}
\]

\subsection{ตัวเก็บประจุต่อขนาน}

ตัวเก็บประจุที่ต่อขนานมีความต่างศักย์เท่ากัน

\[
\Delta V_1=\Delta V_2=\cdots=\Delta V
\]

ประจุรวมที่แหล่งจ่ายส่งให้ชุดคือ

\[
q=q_1+q_2+\cdots
\]

ใช้ $q_i=C_i\Delta V$

\[
q=(C_1+C_2+\cdots)\Delta V
\]

เปรียบเทียบกับ $q=C_{\text{eq}}\Delta V$

\begin{tcolorbox}[title={ตัวเก็บประจุต่อขนาน}]
\[
\boxed{C_{\text{eq}}=C_1+C_2+C_3+\cdots}
\]
\end{tcolorbox}

\subsection{ตัวเก็บประจุต่ออนุกรม}

เมื่อเริ่มจากตัวเก็บประจุไม่มีประจุและต่อเป็นอนุกรม ประจุที่ไหลเข้าสู่แผ่นหนึ่งทำให้แผ่นภายในที่เชื่อมกันเกิดประจุขนาดเท่ากันแต่เครื่องหมายตรงข้าม จึงมีขนาดประจุเท่ากันทุกตัว

\[
q_1=q_2=\cdots=q
\]

ความต่างศักย์รวมคือ

\[
\Delta V=\Delta V_1+\Delta V_2+\cdots
\]

ใช้

\[
\Delta V_i=\frac{q}{C_i}
\]

จึงได้

\[
\Delta V=q\left(\frac{1}{C_1}+\frac{1}{C_2}+\cdots\right)
\]

แต่

\[
\Delta V=\frac{q}{C_{\text{eq}}}
\]

ดังนั้น

\begin{tcolorbox}[title={ตัวเก็บประจุต่ออนุกรม}]
\[
\boxed{\frac{1}{C_{\text{eq}}}
=\frac{1}{C_1}+\frac{1}{C_2}+\frac{1}{C_3}+\cdots}
\]
\end{tcolorbox}

สำหรับสองตัว

\[
\boxed{C_{\text{eq}}=\frac{C_1C_2}{C_1+C_2}}
\]

สังเกตว่ากฎการรวมตัวเก็บประจุตรงข้ามกับตัวต้านทาน

\subsection{อนุมานพลังงานในตัวเก็บประจุ}

ขณะตัวเก็บประจุมีประจุชั่วขณะ $q'$ ความต่างศักย์คือ

\[
V(q')=\frac{q'}{C}
\]

การเพิ่มประจุเล็กน้อย $dq'$ ต้องใช้งาน

\[
dU=V(q')\,dq'
\]

ดังนั้นพลังงานที่ใช้ชาร์จจาก $0$ ถึง $q$ คือ

\[
U=\int_0^q\frac{q'}{C}\,dq'
\]

\[
U=\frac{1}{C}\left[\frac{q'^2}{2}\right]_0^q
\]

\[
U=\frac{q^2}{2C}
\]

ใช้ $q=C\Delta V$

\begin{tcolorbox}[title={พลังงานในตัวเก็บประจุ}]
\[
\boxed{U_C=\frac{q^2}{2C}}
\]

\[
\boxed{U_C=\frac{1}{2}q\Delta V}
\]

\[
\boxed{U_C=\frac{1}{2}C(\Delta V)^2}
\]
\end{tcolorbox}

\subsection{การกระจายประจุใหม่และพลังงานที่เปลี่ยนเป็นความร้อน}

เมื่อตัวเก็บประจุที่มีประจุต่อเข้ากับตัวเก็บประจุอีกตัวผ่านตัวต้านทาน ประจุรวมของระบบตัวนำที่แยกจากภายนอกคงตัว แต่พลังงานในสนามไฟฟ้าไม่จำเป็นต้องคงตัว ส่วนต่างเปลี่ยนเป็นความร้อนในตัวต้านทาน

\[
E_{\text{heat}}=U_i-U_f
\]

ไม่ควรใช้การอนุรักษ์พลังงานในตัวเก็บประจุเพียงอย่างเดียว เพราะตัวต้านทานและการแผ่รังสีรับพลังงานบางส่วน

% =========================================================
\section{พฤติกรรมของตัวเก็บประจุในวงจรกระแสตรง}

\subsection{แรงดันคร่อมตัวเก็บประจุเปลี่ยนฉับพลันไม่ได้}

จาก

\[
i=C\frac{dV_C}{dt}
\]

ซึ่งได้จาก $q=CV_C$ แล้วหาอนุพันธ์ตามเวลา ถ้า $V_C$ เปลี่ยนเป็นค่าจำกัดภายในเวลาศูนย์ จะต้องมีกระแสเป็นอนันต์

ดังนั้นในวงจรจริงที่กระแสมีค่าจำกัด

\[
\boxed{V_C(0^+)=V_C(0^-)}
\]

ถ้าตัวเก็บประจุเริ่มไม่มีประจุ

\[
V_C(0^-)=0
\]

จึงมี

\[
V_C(0^+)=0
\]

ในขณะแรกตัวเก็บประจุจึงมีความต่างศักย์เป็นศูนย์ แต่ข้อความว่า ``ทำตัวเหมือนลวด'' ใช้ได้เฉพาะการวิเคราะห์แรงดันขณะเริ่มต้น ไม่ควรนำไปใช้โดยไม่ตรวจโครงสร้างวงจร

\subsection{สภาวะคงตัวหลังเวลานาน}

สำหรับแหล่งจ่ายกระแสตรง เมื่อเวลานานมาก

\[
\frac{dV_C}{dt}=0
\]

จึงมี

\[
i_C=C\frac{dV_C}{dt}=0
\]

ดังนั้นตัวเก็บประจุทำตัวเป็นวงจรเปิดในสภาวะคงตัว

\begin{tcolorbox}[title={กฎวิเคราะห์ตัวเก็บประจุในวงจร DC}]
\begin{enumerate}
    \item ขณะสับสวิตช์ทันที ใช้ความต่อเนื่องของ $V_C$
    \item หลังเวลานาน ใช้ $i_C=0$ และแทนแขนงตัวเก็บประจุด้วยวงจรเปิด
\end{enumerate}
\end{tcolorbox}

% =========================================================
\section{วงจร RC}

\subsection{การชาร์จตัวเก็บประจุผ่านตัวต้านทาน}

พิจารณาแบตเตอรี่ $\mathcal{E}$ ตัวต้านทาน $R$ และตัวเก็บประจุ $C$ ต่ออนุกรม โดยเริ่มจาก $q(0)=0$

ใช้กฎลูป

\[
\mathcal{E}-iR-\frac{q}{C}=0
\]

และ

\[
i=\frac{dq}{dt}
\]

จึงได้

\[
R\frac{dq}{dt}=\mathcal{E}-\frac{q}{C}
\]

\[
\frac{dq}{dt}=\frac{C\mathcal{E}-q}{RC}
\]

แยกตัวแปร

\[
\frac{dq}{C\mathcal{E}-q}=\frac{dt}{RC}
\]

อินทิเกรตจาก $q'=0$ ที่ $t'=0$ ถึง $q$ ที่ $t$

\[
\int_0^q\frac{dq'}{C\mathcal{E}-q'}
=\int_0^t\frac{dt'}{RC}
\]

\[
-\ln\left(\frac{C\mathcal{E}-q}{C\mathcal{E}}\right)
=\frac{t}{RC}
\]

จัดรูป

\[
C\mathcal{E}-q=C\mathcal{E}e^{-t/(RC)}
\]

ดังนั้น

\begin{tcolorbox}[title={การชาร์จตัวเก็บประจุ}]
\[
\boxed{q(t)=C\mathcal{E}\left(1-e^{-t/(RC)}\right)}
\]

\[
\boxed{V_C(t)=\mathcal{E}\left(1-e^{-t/(RC)}\right)}
\]
\end{tcolorbox}

หากระแสโดยหาอนุพันธ์

\[
i(t)=\frac{dq}{dt}
\]

\[
\boxed{i(t)=\frac{\mathcal{E}}{R}e^{-t/(RC)}}
\]

ตรวจขอบเขต

\[
q(0)=0,\qquad i(0)=\frac{\mathcal{E}}{R}
\]

และเมื่อ $t\to\infty$

\[
q\to C\mathcal{E},\qquad i\to 0
\]

\subsection{การคายประจุผ่านตัวต้านทาน}

ให้ตัวเก็บประจุมีประจุเริ่มต้น $Q_0$ แล้วคายผ่าน $R$

ใช้กฎลูปโดยเลือกทิศกระแสตามทิศประจุบวกไหลออกจากแผ่นบวก

\[
\frac{q}{C}-iR=0
\]

แต่ประจุบนแผ่นบวกลดลง

\[
i=-\frac{dq}{dt}
\]

ดังนั้น

\[
\frac{q}{C}+R\frac{dq}{dt}=0
\]

\[
\frac{dq}{q}=-\frac{dt}{RC}
\]

อินทิเกรตโดยใช้ $q(0)=Q_0$

\[
\ln\left(\frac{q}{Q_0}\right)=-\frac{t}{RC}
\]

จึงได้

\begin{tcolorbox}[title={การคายประจุ}]
\[
\boxed{q(t)=Q_0e^{-t/(RC)}}
\]

\[
\boxed{V_C(t)=V_0e^{-t/(RC)}}
\]

\[
\boxed{|i(t)|=\frac{Q_0}{RC}e^{-t/(RC)}
=\frac{V_0}{R}e^{-t/(RC)}}
\]
\end{tcolorbox}

เครื่องหมายของ $i$ ขึ้นกับทิศที่กำหนด

\subsection{ค่าคงตัวเวลา}

นิยามค่าคงตัวเวลาของวงจร $RC$ เป็น

\[
\boxed{\tau\equiv RC}
\]

ตรวจมิติ

\[
[RC]=\frac{\si{V}}{\si{A}}\frac{\si{C}}{\si{V}}
=\frac{\si{C}}{\si{A}}
=\si{s}
\]

เมื่อ $t=\tau$

สำหรับการชาร์จ

\[
\frac{q}{Q_{\max}}=1-e^{-1}\approx 0.632
\]

สำหรับการคาย

\[
\frac{q}{Q_0}=e^{-1}\approx 0.368
\]

\subsection{กำลังระหว่างการชาร์จ}

จากกฎลูป

\[
\mathcal{E}=iR+\frac{q}{C}
\]

คูณทั้งสมการด้วย $i$

\[
\mathcal{E}i=i^2R+\frac{q}{C}i
\]

โดย

\[
P_{\text{source}}=\mathcal{E}i
\]

\[
P_R=i^2R
\]

และ

\[
\frac{dU_C}{dt}
=\frac{d}{dt}\left(\frac{q^2}{2C}\right)
=\frac{q}{C}\frac{dq}{dt}
=\frac{q}{C}i
\]

ดังนั้น

\begin{tcolorbox}[title={การอนุรักษ์กำลังในวงจร RC}]
\[
\boxed{\mathcal{E}i=i^2R+\frac{dU_C}{dt}}
\]
\end{tcolorbox}

พลังงานจากแหล่งกำเนิดส่วนหนึ่งเพิ่มพลังงานสนามไฟฟ้า และอีกส่วนเปลี่ยนเป็นความร้อนในตัวต้านทาน

% =========================================================
\section{เทคนิควิเคราะห์วงจรซับซ้อน}

\subsection{จุดที่ต่อด้วยลวดอุดมคติมีศักย์เท่ากัน}

ลวดอุดมคติมีความต้านทานเป็นศูนย์ จึงมี

\[
\Delta V=IR=0
\]

ทุกจุดบนโหนดเดียวกันจึงมีศักย์เท่ากัน การวาดวงจรใหม่โดยรักษาการเชื่อมต่อของโหนดไว้ไม่เปลี่ยนวงจรทางไฟฟ้า

\subsection{สมมาตรและจุดศักย์เท่ากัน}

ถ้าวงจรสมมาตรเมื่อสลับจุดสองจุด และเงื่อนไขขอบเขตไม่เปลี่ยน จุดสองจุดนั้นต้องมีศักย์เท่ากัน

ถ้ามีตัวต้านทานเชื่อมระหว่างจุดศักย์เท่ากัน

\[
I=\frac{\Delta V}{R}=0
\]

จึงตัดแขนงนั้นออกได้โดยไม่เปลี่ยนวงจรส่วนอื่น

\subsection{วงจรบันได}

วงจรบันไดควรลดจากด้านปลายย้อนกลับมา หรือกำหนดความต้านทานหรือความจุสมมูลของส่วนที่ซ้ำกันเป็นตัวแปร แล้วตั้งสมการจากโครงสร้างที่ซ้ำ

\subsection{วงจรบริดจ์สมดุล}

สำหรับบริดจ์ตัวต้านทาน ถ้าจุดกึ่งกลางสองจุดมีศักย์เท่ากัน กระแสในแขนงเชื่อมกลางเป็นศูนย์ เงื่อนไขสมดุลมาตรฐานคือ

\[
\boxed{\frac{R_1}{R_2}=\frac{R_3}{R_4}}
\]

สูตรนี้อนุมานได้จากตัวแบ่งแรงดันทั้งสองด้าน ไม่ควรจำโดยไม่เข้าใจว่ามาจากเงื่อนไขศักย์กึ่งกลางเท่ากัน

% =========================================================
\section{วิธีเลือกหลักในการแก้โจทย์}

\begin{tcolorbox}[title={แผนแก้โจทย์วงจร}]
\begin{enumerate}
    \item วาดโหนดใหม่ให้ชัด และทำเครื่องหมายจุดที่มีศักย์เท่ากัน
    \item ถ้าลดอนุกรมหรือขนานได้ ให้ลดก่อน
    \item ถ้ามีตัวเก็บประจุและโจทย์บอก ``ทันที'' ให้ใช้ความต่อเนื่องของ $V_C$
    \item ถ้าโจทย์บอก ``นานมาก'' หรือ ``กระแสคงที่แล้ว'' ให้แทนตัวเก็บประจุด้วยวงจรเปิด
    \item ถ้าลดวงจรไม่ได้ ให้ใช้กฎจุดต่อและกฎลูป
    \item ถ้าถามกำลัง ให้เลือก $P=IV$, $I^2R$ หรือ $V^2/R$ ตามปริมาณที่ทราบและสิ่งที่คงตัว
    \item ถ้าถามการชาร์จหรือคายตามเวลา ให้หา $R_{\text{seen}}$ ที่ตัวเก็บประจุมองเห็น แล้วใช้ $\tau=R_{\text{seen}}C$
    \item ตรวจคำตอบด้วยขอบเขต หน่วย และการอนุรักษ์กำลัง
\end{enumerate}
\end{tcolorbox}

% =========================================================
\section{ชุดโจทย์รวมท้ายบท: ระดับ สอวน. ค่าย 1}

\subsection*{ระดับพื้นฐาน}

\begin{enumerate}
    \item มีกระแสคงตัว $2.0\,\si{A}$ ไหลเป็นเวลา $3.0\,\si{min}$ จงหาประจุที่ผ่านหน้าตัดและจำนวนอิเล็กตรอน โดยใช้ $e=1.60\times10^{-19}\,\si{C}$

    \item ลวดยาว $2.0\,\si{m}$ พื้นที่หน้าตัด $1.0\,\si{mm^2}$ มีสภาพต้านทาน $1.7\times10^{-8}\,\si{\ohm.m}$ จงหาความต้านทาน

    \item ตัวต้านทาน $6\,\si{\ohm}$ ต่อกับความต่างศักย์ $12\,\si{V}$ จงหากระแส กำลัง และพลังงานที่เปลี่ยนเป็นความร้อนในเวลา $5\,\si{min}$

    \item ตัวต้านทาน $3\,\si{\ohm}$ และ $6\,\si{\ohm}$ จงหาความต้านทานสมมูลเมื่อ (ก) ต่ออนุกรม (ข) ต่อขนาน

    \item ตัวเก็บประจุ $4\,\si{\micro\farad}$ ต่อกับแบตเตอรี่ $12\,\si{V}$ จงหาประจุและพลังงานที่เก็บ
\end{enumerate}

\subsection*{ระดับกลาง}

\begin{enumerate}
    \setcounter{enumi}{5}
    \item แบตเตอรี่ $\mathcal{E}=12\,\si{V}$ มีความต้านทานภายใน $1\,\si{\ohm}$ ต่อกับโหลด $5\,\si{\ohm}$ จงหากระแส แรงดันคร่อมขั้ว และกำลังที่โหลดรับ

    \item ตัวต้านทาน $R_1$ ต่ออนุกรมกับวงขนานของ $R_2$ และ $R_3$ จงหาความต้านทานสมมูลในรูปของ $R_1,R_2,R_3$

    \item ตัวเก็บประจุ $C_1$ และ $C_2$ ต่ออนุกรมกับแบตเตอรี่ $V$ จงหาประจุบนแต่ละตัวและความต่างศักย์คร่อมแต่ละตัว

    \item วงจร $RC$ มี $R=2.0\,\si{M\ohm}$ และ $C=5.0\,\si{\micro\farad}$ ต่อกับแบตเตอรี่ $10\,\si{V}$ จงหา $\tau$, $q(\tau)$ และ $i(\tau)$

    \item ตัวเก็บประจุ $C_1$ มีประจุเริ่มต้น $Q_0$ และ $C_2$ ไม่มีประจุ เมื่อต่อขนานขั้วเหมือนกันผ่านตัวต้านทาน จงหาความต่างศักย์สุดท้ายและพลังงานที่เปลี่ยนเป็นความร้อน
\end{enumerate}

\subsection*{ระดับยาก}

\begin{enumerate}
    \setcounter{enumi}{10}
    \item ลวดความต้านทานรวม $R$ ถูกแบ่งเป็นสองส่วน $xR$ และ $(1-x)R$ แล้วต่อขนาน จงหาค่า $x$ ที่ทำให้ความต้านทานสมมูลมากที่สุด และหาค่าสูงสุด

    \item แหล่งกำเนิด $\mathcal{E}$ มีความต้านทานภายใน $r$ ต่อกับโหลด $R$ จงพิสูจน์โดยไม่ใช้แคลคูลัสว่ากำลังที่โหลดรับมากที่สุดเมื่อ $R=r$

    \item วงจรบันไดอนันต์ประกอบด้วยตัวต้านทาน $R$ ต่ออนุกรมหนึ่งตัว แล้วตามด้วยตัวต้านทาน $R$ ลงกราวด์ จากนั้นโครงสร้างซ้ำเดิม จงตั้งสมการหาความต้านทานสมมูล

    \item ตัวเก็บประจุ $C$ เริ่มมีแรงดัน $V_0$ และคายประจุผ่าน $R$ จงหาเวลาที่พลังงานในตัวเก็บประจุเหลือ $1/4$ ของค่าเริ่มต้น
\end{enumerate}

% =========================================================
\section{เฉลยย่อชุดโจทย์ท้ายบท}

\begin{enumerate}
    \item
    \[
    \Delta q=It=(2.0)(180)=360\,\si{C}
    \]
    \[
    N=\frac{360}{1.60\times10^{-19}}
    =2.25\times10^{21}
    \]

    \item
    \[
    R=\rho\frac{L}{A}
    =\frac{(1.7\times10^{-8})(2.0)}{1.0\times10^{-6}}
    =3.4\times10^{-2}\,\si{\ohm}
    \]

    \item
    \[
    I=\frac{12}{6}=2\,\si{A}
    \]
    \[
    P=IV=24\,\si{W}
    \]
    \[
    E=Pt=(24)(300)=7.2\times10^3\,\si{J}
    \]

    \item
    \[
    R_{\text{series}}=9\,\si{\ohm}
    \]
    \[
    R_{\text{parallel}}=\frac{(3)(6)}{3+6}=2\,\si{\ohm}
    \]

    \item
    \[
    q=CV=(4\,\si{\micro\farad})(12\,\si{V})
    =48\,\si{\micro\coulomb}
    \]
    \[
    U=\frac12CV^2=288\,\si{\micro\joule}
    \]

    \item
    \[
    I=\frac{12}{5+1}=2\,\si{A}
    \]
    \[
    V_{\text{terminal}}=\mathcal{E}-Ir=10\,\si{V}
    \]
    \[
    P_R=I^2R=20\,\si{W}
    \]

    \item
    \[
    R_{\text{eq}}=R_1+\frac{R_2R_3}{R_2+R_3}
    \]

    \item
    \[
    C_{\text{eq}}=\frac{C_1C_2}{C_1+C_2}
    \]
    \[
    q_1=q_2=C_{\text{eq}}V
    \]
    \[
    V_1=\frac{q}{C_1}=V\frac{C_2}{C_1+C_2}
    \]
    \[
    V_2=\frac{q}{C_2}=V\frac{C_1}{C_1+C_2}
    \]

    \item
    \[
    \tau=RC=(2.0\times10^6)(5.0\times10^{-6})=10\,\si{s}
    \]
    \[
    q(\tau)=C\mathcal{E}(1-e^{-1})
    \approx31.6\,\si{\micro\coulomb}
    \]
    \[
    i(\tau)=\frac{\mathcal{E}}{R}e^{-1}
    \approx1.84\,\si{\micro\ampere}
    \]

    \item
    \[
    V_f=\frac{Q_0}{C_1+C_2}
    \]
    \[
    E_{\text{heat}}
    =\frac{Q_0^2}{2C_1}
    -\frac{Q_0^2}{2(C_1+C_2)}
    =\frac{Q_0^2C_2}{2C_1(C_1+C_2)}
    \]

    \item
    \[
    R_{\parallel}=\frac{xR(1-x)R}{R}=x(1-x)R
    \]
    \[
    x(1-x)=\frac14-\left(x-\frac12\right)^2
    \]
    จึงมากที่สุดเมื่อ
    \[
    x=\frac12
    \]
    และ
    \[
    R_{\parallel,\max}=\frac{R}{4}
    \]

    \item
    \[
    P_R=\frac{\mathcal{E}^2R}{(R+r)^2}
    \]
    เพราะ
    \[
    (R+r)^2\geq4Rr
    \]
    จึงได้
    \[
    P_R\leq\frac{\mathcal{E}^2}{4r}
    \]
    โดยเท่ากันเมื่อ $R=r$

    \item ให้ความต้านทานสมมูลของบันไดเป็น $X$ ส่วนที่เหลือหลังตัวต้านทานอนุกรมตัวแรกยังมีค่า $X$ จึงได้
    \[
    X=R+(R\parallel X)
    \]
    \[
    X=R+\frac{RX}{R+X}
    \]
    จัดรูป
    \[
    X^2-RX-R^2=0
    \]
    รากบวกคือ
    \[
    \boxed{X=\frac{1+\sqrt5}{2}R}
    \]

    \item พลังงานแปรตาม $V_C^2$
    \[
    \frac{U}{U_0}=e^{-2t/(RC)}
    \]
    ให้ $U/U_0=1/4$
    \[
    e^{-2t/(RC)}=\frac14
    \]
    \[
    \boxed{t=RC\ln2}
    \]
\end{enumerate}

% =========================================================
\section{สรุปสูตรสำคัญ}

\begin{tcolorbox}[title={สรุป: กระแส ความต้านทาน และกำลัง}]
\[
I=\frac{dq}{dt}
\]

\[
I=nq_cAv_d
\]

\[
R=\frac{\Delta V}{I}
\]

\[
R=\rho\frac{L}{A}
\]

\[
\sigma=\frac{1}{\rho}
\]

\[
P=I\Delta V=I^2R=\frac{(\Delta V)^2}{R}
\]

\[
V_{\text{terminal}}=\mathcal{E}-Ir
\]
\end{tcolorbox}

\begin{tcolorbox}[title={สรุป: การรวมวงจรและตัวเก็บประจุ}]
ตัวต้านทานอนุกรม

\[
R_{\text{eq}}=\sum_iR_i
\]

ตัวต้านทานขนาน

\[
\frac{1}{R_{\text{eq}}}=\sum_i\frac{1}{R_i}
\]

ตัวเก็บประจุขนาน

\[
C_{\text{eq}}=\sum_iC_i
\]

ตัวเก็บประจุอนุกรม

\[
\frac{1}{C_{\text{eq}}}=\sum_i\frac{1}{C_i}
\]

\[
q=C\Delta V
\]

\[
U_C=\frac{q^2}{2C}=\frac12q\Delta V=\frac12C(\Delta V)^2
\]
\end{tcolorbox}

\begin{tcolorbox}[title={สรุป: วงจร RC}]
\[
\tau=RC
\]

การชาร์จ

\[
q(t)=C\mathcal{E}\left(1-e^{-t/(RC)}\right)
\]

\[
i(t)=\frac{\mathcal{E}}{R}e^{-t/(RC)}
\]

การคาย

\[
q(t)=Q_0e^{-t/(RC)}
\]

\[
|i(t)|=\frac{Q_0}{RC}e^{-t/(RC)}
\]

การอนุรักษ์กำลัง

\[
\mathcal{E}i=i^2R+\frac{dU_C}{dt}
\]
\end{tcolorbox}

% =========================================================
\section{ข้อสอบจริง สอวน. ที่ใช้ความรู้บทวงจรไฟฟ้า}

\begin{tcolorbox}[title={ขอบเขตของข้อสอบท้ายบท}]
ชุดนี้รวมโจทย์ที่ใช้ความรู้เรื่องกระแส ความต้านทาน สภาพต้านทาน กำลังไฟฟ้า แรงเคลื่อนไฟฟ้า กฎของเคอร์ชอฟฟ์ การต่อตัวต้านทาน ตัวเก็บประจุ และวงจร $RC$

โจทย์แรงระหว่างแผ่นตัวเก็บประจุซึ่งเป็นไฟฟ้าสถิตโดยตรงควรอยู่ในบทไฟฟ้าสถิต จึงไม่ทำซ้ำในบทนี้
\end{tcolorbox}

\begin{tcolorbox}[title={คำแนะนำการใส่ภาพข้อสอบ}]
ให้ตัดภาพข้อสอบจริงแล้วบันทึกตามชื่อไฟล์ที่กำหนด ถ้ายังไม่มีไฟล์ ภาพจะไม่แสดง แต่เอกสารยังคอมไพล์ได้
\end{tcolorbox}

% =========================================================
\subsection{ปี 2560 ข้อ 13: ความต้านทานไม่ทราบค่าในวงจรขนาน}

\vspace{1em}

\IfFileExists{img/posn60q13.png}{
\begin{center}
    \includegraphics[width=1\linewidth]{img/posn60q13.png}
\end{center}
}{
\begin{tcolorbox}[title=ตำแหน่งภาพที่แนะนำ]
ให้ตัดภาพข้อสอบปี 2560 ข้อ 13 แล้วบันทึกเป็น
\[
\texttt{img/posn60q13.png}
\]
\end{tcolorbox}
}

\begin{examplebox}{เฉลยละเอียด}
ตัวต้านทาน $X$, $15\,\si{\ohm}$, $6\,\si{\ohm}$ และ $10\,\si{\ohm}$ ต่ออยู่ระหว่างจุด $a$ และ $b$ เดียวกัน จึงต่อขนานทั้งหมด

โอห์มมิเตอร์อ่านความต้านทานสมมูล

\[
R_{\text{eq}}=2.0\,\si{\ohm}
\]

ดังนั้น

\[
\frac{1}{2}
=\frac{1}{X}+\frac{1}{15}+\frac{1}{6}+\frac{1}{10}
\]

รวมสามพจน์ที่ทราบ

\[
\frac{1}{15}+\frac{1}{6}+\frac{1}{10}
=\frac{2+5+3}{30}
=\frac{1}{3}
\]

จึงได้

\[
\frac{1}{2}=\frac{1}{X}+\frac{1}{3}
\]

\[
\frac{1}{X}=\frac{1}{6}
\]

\[
\boxed{X=6.0\,\si{\ohm}}
\]

ตอบ D
\end{examplebox}

% ---------------------------------------------------------
\subsection{ปี 2560 ข้อ 14: วงจรตัวต้านทานและตัวเก็บประจุในสภาวะคงตัว}

\vspace{1em}

\IfFileExists{img/posn60q14.png}{
\begin{center}
    \includegraphics[width=1\linewidth]{img/posn60q14.png}
\end{center}
}{
\begin{tcolorbox}[title=ตำแหน่งภาพที่แนะนำ]
ให้ตัดภาพข้อสอบปี 2560 ข้อ 14 แล้วบันทึกเป็น
\[
\texttt{img/posn60q14.png}
\]
\end{tcolorbox}
}

\begin{examplebox}{เฉลยละเอียด}
เมื่อปล่อยไว้นานจนตัวเก็บประจุเต็ม กระแสผ่านตัวเก็บประจุทุกแขนงเป็นศูนย์ จึงวิเคราะห์แรงดันด้วยวงจรตัวต้านทานก่อน

ตัวต้านทาน $3\,\si{\ohm}$ และ $6\,\si{\ohm}$ ต่อขนาน

\[
R_L=\frac{(3)(6)}{3+6}=2\,\si{\ohm}
\]

ตัวต้านทาน $2.5\,\si{\ohm}$ และ $1.5\,\si{\ohm}$ ต่ออนุกรม

\[
R_R=2.5+1.5=4\,\si{\ohm}
\]

สองส่วนนี้ต่ออนุกรมกับแหล่งจ่าย $12\,\si{V}$

\[
I=\frac{12}{2+4}=2\,\si{A}
\]

แรงดันคร่อมส่วนด้านขวา ซึ่งเป็นแรงดันคร่อมแขนงตัวเก็บประจุ คือ

\[
V_R=IR_R=(2)(4)=8\,\si{V}
\]

ตัวเก็บประจุ $1\,\si{\micro\farad}$ และ $4\,\si{\micro\farad}$ ต่ออนุกรม

\[
C_s=\frac{(1)(4)}{1+4}
=0.8\,\si{\micro\farad}
\]

ประจุในชุดอนุกรมมีขนาดเท่ากันทุกตัว

\[
q=C_sV_R
\]

\[
q=(0.8)(8)=6.4\,\si{\micro\coulomb}
\]

ดังนั้นตัวเก็บประจุ $4\,\si{\micro\farad}$ มีประจุ

\[
\boxed{q=6.4\,\si{\micro\coulomb}}
\]

ตอบ A
\end{examplebox}

% ---------------------------------------------------------
\subsection{ปี 2560 ข้อ 15: ลวดถูกดึงและกำลังที่ความต่างศักย์คงตัว}

\vspace{1em}

\IfFileExists{img/posn60q15.png}{
\begin{center}
    \includegraphics[width=1\linewidth]{img/posn60q15.png}
\end{center}
}{
\begin{tcolorbox}[title=ตำแหน่งภาพที่แนะนำ]
ให้ตัดภาพข้อสอบปี 2560 ข้อ 15 แล้วบันทึกเป็น
\[
\texttt{img/posn60q15.png}
\]
\end{tcolorbox}
}

\begin{examplebox}{เฉลยละเอียด}
ลวดเดิมมีความยาว $L$ และพื้นที่หน้าตัด $A$

เมื่อความยาวใหม่เป็น

\[
L'=xL
\]

และปริมาตรคงตัว

\[
AL=A'L'
\]

จึงได้

\[
A'=\frac{A}{x}
\]

ความต้านทานใหม่คือ

\[
R'=\rho\frac{L'}{A'}
=\rho\frac{xL}{A/x}
=x^2R
\]

แหล่งจ่ายรักษาความต่างศักย์คงตัว จึงใช้

\[
P=\frac{V^2}{R}
\]

ดังนั้น

\[
\frac{P'}{P}
=\frac{V^2/(x^2R)}{V^2/R}
\]

\[
\boxed{\frac{P'}{P}=\frac{1}{x^2}}
\]

ตอบ D
\end{examplebox}

% ---------------------------------------------------------
\subsection{ปี 2560 ตอนที่ 2 ข้อ 7: วงจรสามแขนงและกฎของเคอร์ชอฟฟ์}

\vspace{1em}

\IfFileExists{img/posn60p2q7.png}{
\begin{center}
    \includegraphics[width=1\linewidth]{img/posn60p2q7.png}
\end{center}
}{
\begin{tcolorbox}[title=ตำแหน่งภาพที่แนะนำ]
ให้ตัดภาพข้อสอบปี 2560 ตอนที่ 2 ข้อ 7 แล้วบันทึกเป็น
\[
\texttt{img/posn60p2q7.png}
\]
\end{tcolorbox}
}

\begin{examplebox}{เฉลยละเอียด}
ให้

\[
V\equiv V_{\text{left}}-V_{\text{right}}
\]

และกำหนดกระแสทางขวาเป็นบวกในทุกแขนง

แขนงบนมีกระแสให้มา

\[
I_1=2.0\,\si{A}
\]

เมื่อเดินจากซ้ายไปขวา แบตเตอรี่ $3.0\,\si{V}$ ทำให้ศักย์เพิ่มขึ้น และตัวต้านทาน $1.0\,\si{\ohm}$ ทำให้ศักย์ลดลง จึงมี

\[
V=I_1(1.0)-3.0
\]

\[
V=(2.0)(1.0)-3.0=-1.0\,\si{V}
\]

ให้ $J_2$ เป็นกระแสแขนงกลางในทิศซ้ายไปขวา แบตเตอรี่ $6.0\,\si{V}$ มีขั้วบวกทางซ้าย จึงได้

\[
V=6.0+2.0J_2
\]

แทน $V=-1.0\,\si{V}$

\[
-1.0=6.0+2.0J_2
\]

\[
J_2=-3.5\,\si{A}
\]

เครื่องหมายลบแสดงว่ากระแสจริงไหลไปทางซ้ายด้วยขนาด $3.5\,\si{A}$

ใช้กฎจุดต่อ

\[
I_1+J_2+I_3=0
\]

\[
2.0-3.5+I_3=0
\]

ดังนั้น

\[
\boxed{I_3=1.5\,\si{A}}
\]

ทิศไปทางขวาตามลูกศรในรูป
\end{examplebox}

% =========================================================
\subsection{ปี 2561 ข้อ 13: ค่าสูงสุดและต่ำสุดจากตัวต้านทานสามตัว}

\vspace{1em}

\IfFileExists{img/posn61q13.png}{
\begin{center}
    \includegraphics[width=1\linewidth]{img/posn61q13.png}
\end{center}
}{
\begin{tcolorbox}[title=ตำแหน่งภาพที่แนะนำ]
ให้ตัดภาพข้อสอบปี 2561 ข้อ 13 แล้วบันทึกเป็น
\[
\texttt{img/posn61q13.png}
\]
\end{tcolorbox}
}

\begin{examplebox}{เฉลยละเอียด}
ตัวต้านทานมีค่า $R$, $R$ และ $2R$

ค่ามากที่สุดได้จากการต่ออนุกรม

\[
R_{\max}=R+R+2R=4R
\]

ค่าน้อยที่สุดได้จากการต่อขนาน

\[
\frac{1}{R_{\min}}
=\frac{1}{R}+\frac{1}{R}+\frac{1}{2R}
=\frac{5}{2R}
\]

\[
R_{\min}=\frac{2R}{5}
\]

ดังนั้น

\[
\frac{R_{\max}}{R_{\min}}
=\frac{4R}{2R/5}
=10
\]

\[
\boxed{\frac{R_{\max}}{R_{\min}}=10}
\]

ตอบ D
\end{examplebox}

% ---------------------------------------------------------
\subsection{ปี 2561 ข้อ 14: ตัวเก็บประจุในวงจรผสม}

\vspace{1em}

\IfFileExists{img/posn61q14.png}{
\begin{center}
    \includegraphics[width=1\linewidth]{img/posn61q14.png}
\end{center}
}{
\begin{tcolorbox}[title=ตำแหน่งภาพที่แนะนำ]
ให้ตัดภาพข้อสอบปี 2561 ข้อ 14 แล้วบันทึกเป็น
\[
\texttt{img/posn61q14.png}
\]
\end{tcolorbox}
}

\begin{examplebox}{เฉลยละเอียด}
วงจรมีสองแขนงขนานต่อกับแหล่งจ่าย $48.0\,\si{V}$

พิจารณาแขนงที่มีตัวเก็บประจุ $6.00\,\si{\micro\farad}$

ตัวเก็บประจุ $2.00$, $3.00$ และ $7.00\,\si{\micro\farad}$ ต่อขนาน

\[
C_p=2+3+7=12\,\si{\micro\farad}
\]

ชุดนี้ต่ออนุกรมกับ $6.00\,\si{\micro\farad}$

\[
C_s=\frac{(6)(12)}{6+12}
=4.00\,\si{\micro\farad}
\]

ทั้งแขนงมีแรงดัน $48.0\,\si{V}$ จึงมีประจุในชุดอนุกรม

\[
q=C_sV
\]

\[
q=(4.00)(48.0)
=192\,\si{\micro\coulomb}
\]

ประจุบนตัวเก็บประจุทุกส่วนในอนุกรมมีขนาดเท่ากัน ดังนั้นตัว $6.00\,\si{\micro\farad}$ มี

\[
\boxed{q=192\,\si{\micro\coulomb}}
\]

ตอบ C
\end{examplebox}

% ---------------------------------------------------------
\subsection{ปี 2561 ข้อ 15: หารัศมีลวดจากกราฟ $I$ กับ $V$}

\vspace{1em}

\IfFileExists{img/posn61q15.png}{
\begin{center}
    \includegraphics[width=1\linewidth]{img/posn61q15.png}
\end{center}
}{
\begin{tcolorbox}[title=ตำแหน่งภาพที่แนะนำ]
ให้ตัดภาพข้อสอบปี 2561 ข้อ 15 แล้วบันทึกเป็น
\[
\texttt{img/posn61q15.png}
\]
\end{tcolorbox}
}

\begin{examplebox}{เฉลยละเอียด}
จากกราฟ เมื่อ

\[
V=30\,\si{V}
\]

มีกระแส

\[
I=60\,\si{A}
\]

ดังนั้นความต้านทานคือ

\[
R=\frac{V}{I}=\frac{30}{60}=0.50\,\si{\ohm}
\]

โจทย์ให้สภาพนำไฟฟ้า

\[
\sigma=\frac{2.5\times10^6}{\pi}\,\si{(\ohm.m)^{-1}}
\]

จึงมีสภาพต้านทาน

\[
\rho=\frac{1}{\sigma}
=\frac{\pi}{2.5\times10^6}\,\si{\ohm.m}
\]

ใช้

\[
R=\rho\frac{L}{A}
\]

โดย $L=0.20\,\si{m}$ และ $A=\pi r^2$

\[
0.50
=\frac{\pi}{2.5\times10^6}
\frac{0.20}{\pi r^2}
\]

ตัด $\pi$

\[
r^2=1.6\times10^{-7}\,\si{m^2}
\]

\[
r=4.0\times10^{-4}\,\si{m}
\]

\[
\boxed{r=0.40\,\si{mm}}
\]

ตอบ C
\end{examplebox}

% ---------------------------------------------------------
\subsection{ปี 2561 ตอนที่ 2 ข้อ 7: การเปลี่ยนค่าที่อ่านได้จากแอมมิเตอร์}

\vspace{1em}

\IfFileExists{img/posn61p2q7.png}{
\begin{center}
    \includegraphics[width=1\linewidth]{img/posn61p2q7.png}
\end{center}
}{
\begin{tcolorbox}[title=ตำแหน่งภาพที่แนะนำ]
ให้ตัดภาพข้อสอบปี 2561 ตอนที่ 2 ข้อ 7 แล้วบันทึกเป็น
\[
\texttt{img/posn61p2q7.png}
\]
\end{tcolorbox}
}

\begin{examplebox}{เฉลยละเอียด}
\textbf{ก่อนปิดสวิตช์}

แขนงตัวต้านทาน $20\,\si{\ohm}$ ทั้งสองยังไม่ต่อครบวงจร เหลือเพียง $15\,\si{\ohm}$ อนุกรมกับ $10\,\si{\ohm}$

\[
I_{A,\text{open}}
=\frac{50}{15+10}
=2.00\,\si{A}
\]

\textbf{หลังปิดสวิตช์}

ตัวต้านทาน $10\,\si{\ohm}$, $20\,\si{\ohm}$ และ $20\,\si{\ohm}$ ต่อขนาน

\[
\frac{1}{R_p}
=\frac{1}{10}+\frac{1}{20}+\frac{1}{20}
=\frac{1}{5}
\]

\[
R_p=5\,\si{\ohm}
\]

ต่ออนุกรมกับ $15\,\si{\ohm}$

\[
I_{\text{total}}=\frac{50}{15+5}=2.50\,\si{A}
\]

แรงดันคร่อมชุดขนานคือ

\[
V_p=I_{\text{total}}R_p=(2.50)(5)=12.5\,\si{V}
\]

กระแสผ่านแขนง $10\,\si{\ohm}$ ที่แอมมิเตอร์วัดคือ

\[
I_{A,\text{closed}}=\frac{12.5}{10}=1.25\,\si{A}
\]

ดังนั้นค่าที่อ่านได้ลดลง

\[
\Delta I=2.00-1.25
\]

\[
\boxed{\Delta I=0.75\,\si{A}}
\]
\end{examplebox}

% =========================================================
\subsection{ปี 2562 ข้อ 13: กระแสผ่านตัวต้านทานกลางในวงจรสมมาตร}

\vspace{1em}

\IfFileExists{img/posn62q13.png}{
\begin{center}
    \includegraphics[width=1\linewidth]{img/posn62q13.png}
\end{center}
}{
\begin{tcolorbox}[title=ตำแหน่งภาพที่แนะนำ]
ให้ตัดภาพข้อสอบปี 2562 ข้อ 13 แล้วบันทึกเป็น
\[
\texttt{img/posn62q13.png}
\]
\end{tcolorbox}
}

\begin{examplebox}{เฉลยละเอียด}
ให้ความต่างศักย์ระหว่างปลายบนและปลายล่างของตัวต้านทาน $r$ เป็น

\[
V
\]

กำหนดกระแสจากจุดบนลงจุดล่างเป็นบวก

กระแสผ่านตัวต้านทานกลางคือ

\[
I_r=\frac{V}{r}
\]

สำหรับแขนงซ้ายและขวา กระแสในทิศจากบนลงล่างเป็น

\[
I_L=\frac{V-\mathcal{E}}{R}
\]

และ

\[
I_R=\frac{V+\mathcal{E}}{R}
\]

ใช้กฎจุดต่อที่จุดบน

\[
I_r+I_L+I_R=0
\]

\[
\frac{V}{r}
+\frac{V-\mathcal{E}}{R}
+\frac{V+\mathcal{E}}{R}=0
\]

พจน์แรงเคลื่อนไฟฟ้าหักล้างกัน

\[
V\left(\frac{1}{r}+\frac{2}{R}\right)=0
\]

ดังนั้น

\[
V=0
\]

จึงได้

\[
\boxed{I_r=0}
\]

ตอบข้อ 1
\end{examplebox}

% ---------------------------------------------------------
\subsection{ปี 2562 ข้อ 14: กระแสทันทีหลังย้ายสวิตช์}

\vspace{1em}

\IfFileExists{img/posn62q14.png}{
\begin{center}
    \includegraphics[width=1\linewidth]{img/posn62q14.png}
\end{center}
}{
\begin{tcolorbox}[title=ตำแหน่งภาพที่แนะนำ]
ให้ตัดภาพข้อสอบปี 2562 ข้อ 14 แล้วบันทึกเป็น
\[
\texttt{img/posn62q14.png}
\]
\end{tcolorbox}
}

\begin{examplebox}{เฉลยละเอียด}
ก่อนย้ายสวิตช์ ตัวเก็บประจุต่อที่ตำแหน่ง $(1)$ เป็นเวลานาน จึงมีแรงดันเท่ากับแรงดันคร่อม $R_2$ ในตัวแบ่งแรงดัน

\[
V_C(0^-)
=\mathcal{E}\frac{R_2}{R_1+R_2}
\]

แรงดันคร่อมตัวเก็บประจุเปลี่ยนฉับพลันไม่ได้

\[
V_C(0^+)=V_C(0^-)
\]

ทันทีหลังย้ายสวิตช์ไปตำแหน่ง $(2)$ ตัวเก็บประจุคายผ่านตัวต้านทาน $R$

\[
I(0^+)=\frac{V_C(0^+)}{R}
\]

ดังนั้น

\[
\boxed{
I(0^+)
=\frac{R_2}{R_1+R_2}\frac{\mathcal{E}}{R}
}
\]

ตอบข้อ 4
\end{examplebox}

% ---------------------------------------------------------
\subsection{ปี 2562 ข้อ 17: กำลังสูงสุดที่โหลดรับ}

\vspace{1em}

\IfFileExists{img/posn62q17.png}{
\begin{center}
    \includegraphics[width=1\linewidth]{img/posn62q17.png}
\end{center}
}{
\begin{tcolorbox}[title=ตำแหน่งภาพที่แนะนำ]
ให้ตัดภาพข้อสอบปี 2562 ข้อ 17 แล้วบันทึกเป็น
\[
\texttt{img/posn62q17.png}
\]
\end{tcolorbox}
}

\begin{examplebox}{เฉลยละเอียด}
กระแสในวงจรคือ

\[
i=\frac{\mathcal{E}}{r+x}
\]

กำลังที่ตัวต้านทานโหลด $x$ รับคือ

\[
\beta=i^2x
\]

แทนกระแส

\[
\beta=\frac{\mathcal{E}^2x}{(r+x)^2}
\]

ใช้

\[
(r+x)^2\geq4rx
\]

จะได้

\[
\beta
\leq\frac{\mathcal{E}^2x}{4rx}
\]

\[
\beta\leq\frac{\mathcal{E}^2}{4r}
\]

เครื่องหมายเท่ากับเกิดเมื่อ

\[
x=r
\]

ดังนั้นค่ามากที่สุดที่เป็นไปได้คือ

\[
\boxed{\beta_{\max}=\frac{\mathcal{E}^2}{4r}}
\]
\end{examplebox}

% ---------------------------------------------------------
\subsection{ปี 2562 ข้อ 21: หาค่าตัวเก็บประจุจากเงื่อนไขความจุสมมูล}

\vspace{1em}

\IfFileExists{img/posn62q21.png}{
\begin{center}
    \includegraphics[width=1\linewidth]{img/posn62q21.png}
\end{center}
}{
\begin{tcolorbox}[title=ตำแหน่งภาพที่แนะนำ]
ให้ตัดภาพข้อสอบปี 2562 ข้อ 21 แล้วบันทึกเป็น
\[
\texttt{img/posn62q21.png}
\]
\end{tcolorbox}
}

\begin{examplebox}{เฉลยละเอียด}
ตัวเก็บประจุ $C_2$ และ $C$ ต่อขนาน

\[
C_p=C_2+C
\]

ชุดนี้ต่ออนุกรมกับ $C_1$

\[
C_{AB}=\frac{C_1(C_2+C)}{C_1+C_2+C}
\]

โจทย์กำหนดว่า

\[
C_{AB}=C
\]

จึงได้

\[
C(C_1+C_2+C)=C_1(C_2+C)
\]

ขยายและตัดพจน์ $CC_1$

\[
C^2+C_2C-C_1C_2=0
\]

ใช้สูตรสมการกำลังสอง

\[
C=\frac{-C_2\pm\sqrt{C_2^2+4C_1C_2}}{2}
\]

ความจุต้องเป็นบวก จึงเลือกรากบวก

\[
\boxed{
C=\frac{\sqrt{C_2^2+4C_1C_2}-C_2}{2}
}
\]
\end{examplebox}

% =========================================================
\subsection{ปี 2564 ข้อ 10: พลังงานที่เปลี่ยนเป็นความร้อนเมื่อประจุกระจายใหม่}

\vspace{1em}

\IfFileExists{img/posn64q10.png}{
\begin{center}
    \includegraphics[width=1\linewidth]{img/posn64q10.png}
\end{center}
}{
\begin{tcolorbox}[title=ตำแหน่งภาพที่แนะนำ]
ให้ตัดภาพข้อสอบปี 2564 ข้อ 10 แล้วบันทึกเป็น
\[
\texttt{img/posn64q10.png}
\]
\end{tcolorbox}
}

\begin{examplebox}{เฉลยละเอียด}
ตัวเก็บประจุ $C_1$ มีประจุเริ่มต้น $q_0$ ส่วน $C_2$ ไม่มีประจุ

หลังปิดสวิตช์ ประจุรวมคงตัว

\[
q_{\text{total}}=q_0
\]

ตัวเก็บประจุทั้งสองมีแรงดันสุดท้ายเท่ากัน

\[
V_f=\frac{q_0}{C_1+C_2}
\]

พลังงานเริ่มต้นคือ

\[
U_i=\frac{q_0^2}{2C_1}
\]

พลังงานสุดท้ายคือ

\[
U_f=\frac12(C_1+C_2)V_f^2
\]

\[
U_f
=\frac12(C_1+C_2)
\left(\frac{q_0}{C_1+C_2}\right)^2
\]

\[
U_f=\frac{q_0^2}{2(C_1+C_2)}
\]

พลังงานที่เปลี่ยนเป็นความร้อนในตัวต้านทานคือ

\[
Q_R=U_i-U_f
\]

\[
Q_R
=\frac{q_0^2}{2}
\left(\frac{1}{C_1}-\frac{1}{C_1+C_2}\right)
\]

ดังนั้น

\[
\boxed{
Q_R=\frac{C_2q_0^2}{2C_1(C_1+C_2)}
}
\]

ตอบ ค
\end{examplebox}

% ---------------------------------------------------------
\subsection{ปี 2564 ข้อ 11: ตัวเก็บประจุเชื่อมจุดกึ่งกลางของตัวแบ่งแรงดัน}

\vspace{1em}

\IfFileExists{img/posn64q11.png}{
\begin{center}
    \includegraphics[width=1\linewidth]{img/posn64q11.png}
\end{center}
}{
\begin{tcolorbox}[title=ตำแหน่งภาพที่แนะนำ]
ให้ตัดภาพข้อสอบปี 2564 ข้อ 11 แล้วบันทึกเป็น
\[
\texttt{img/posn64q11.png}
\]
\end{tcolorbox}
}

\begin{examplebox}{เฉลยละเอียด}
ในสภาวะคงตัว ตัวเก็บประจุไม่มีกระแสผ่าน จึงหาแรงดันที่จุดกึ่งกลางจากตัวแบ่งแรงดันสองแขนง

แขนงบนมี $R$ และ $2R$ ต่ออนุกรม กระแสคือ

\[
I_t=\frac{\mathcal{E}}{3R}
\]

แรงดันตกจากด้านซ้ายถึงจุดกึ่งกลางบนคือ

\[
V_t=I_tR=\frac{\mathcal{E}}{3}
\]

แขนงล่างมี $3R$ และ $4R$ ต่ออนุกรม กระแสคือ

\[
I_b=\frac{\mathcal{E}}{7R}
\]

แรงดันตกจากด้านซ้ายถึงจุดกึ่งกลางล่างคือ

\[
V_b=I_b(3R)=\frac{3\mathcal{E}}{7}
\]

ขนาดความต่างศักย์คร่อมตัวเก็บประจุคือ

\[
\Delta V_C
=\left|\frac{3\mathcal{E}}{7}-\frac{\mathcal{E}}{3}\right|
\]

\[
\Delta V_C=\frac{2\mathcal{E}}{21}
\]

ดังนั้น

\[
q=C\Delta V_C
\]

\[
\boxed{q=\frac{2}{21}C\mathcal{E}}
\]

ตอบ ก
\end{examplebox}

% ---------------------------------------------------------
\subsection{ปี 2564 ข้อ 20: วงจรบันไดตัวต้านทาน}

\vspace{1em}

\IfFileExists{img/posn64q20.png}{
\begin{center}
    \includegraphics[width=1\linewidth]{img/posn64q20.png}
\end{center}
}{
\begin{tcolorbox}[title=ตำแหน่งภาพที่แนะนำ]
ให้ตัดภาพข้อสอบปี 2564 ข้อ 20 แล้วบันทึกเป็น
\[
\texttt{img/posn64q20.png}
\]
\end{tcolorbox}
}

\begin{examplebox}{เฉลยละเอียด}
ต้องการให้ความต้านทานรวมระหว่าง $A$ และ $B$ เท่ากับ

\[
112\,\si{\ohm}
\]

ตัวต้านทาน $56\,\si{\ohm}$ ตัวแรกต่ออนุกรมกับส่วนที่เหลือ ดังนั้นส่วนที่เหลือต้องมีความต้านทาน

\[
112-56=56\,\si{\ohm}
\]

ที่จุดต่อแรก มีแขนง $112\,\si{\ohm}$ ขนานกับแขนงขวา ให้ความต้านทานแขนงขวาเป็น $Y$

\[
112\parallel Y=56
\]

เพราะตัวต้านทานสองตัวที่เท่ากันต่อขนานให้ครึ่งหนึ่ง จึงได้

\[
Y=112\,\si{\ohm}
\]

แขนงขวามี $56\,\si{\ohm}$ อนุกรมกับส่วนถัดไป จึงต้องมีส่วนถัดไป

\[
X=112-56=56\,\si{\ohm}
\]

ที่จุดต่อที่สอง

\[
X=112\parallel(65+R)
\]

ต้องการ $X=56\,\si{\ohm}$ จึงต้องมี

\[
65+R=112
\]

ดังนั้น

\[
\boxed{R=47\,\si{\ohm}}
\]

ตอบ ค
\end{examplebox}

% ---------------------------------------------------------
\subsection{ปี 2564 ตอนที่ 2 ข้อ 26: อัตราส่วนกระแสในวงจรสองแหล่งกำเนิด}

\vspace{1em}

\IfFileExists{img/posn64p2q26.png}{
\begin{center}
    \includegraphics[width=1\linewidth]{img/posn64p2q26.png}
\end{center}
}{
\begin{tcolorbox}[title=ตำแหน่งภาพที่แนะนำ]
ให้ตัดภาพข้อสอบปี 2564 ตอนที่ 2 ข้อ 26 แล้วบันทึกเป็น
\[
\texttt{img/posn64p2q26.png}
\]
\end{tcolorbox}
}

\begin{examplebox}{เฉลยละเอียด}
เลือกสายล่างเป็นศักย์อ้างอิง

แหล่งกำเนิด $\mathcal{E}_2$ ต่อคร่อมระหว่างสายบนด้านขวากับสายล่างโดยตรง จึงกำหนดศักย์สายบนด้านขวาเป็น

\[
V=\mathcal{E}_2
\]

กระแสผ่าน $R_2$ คือ

\[
i_2=\frac{\mathcal{E}_2}{R_2}
\]

ด้านซ้ายของ $R_1$ มีศักย์ $\mathcal{E}_1$ จึงได้

\[
i_1=\frac{\mathcal{E}_1-\mathcal{E}_2}{R_1}
\]

ดังนั้น

\[
\frac{i_1}{i_2}
=\frac{(\mathcal{E}_1-\mathcal{E}_2)/R_1}
{\mathcal{E}_2/R_2}
\]

\[
\boxed{
\frac{i_1}{i_2}
=\frac{R_2(\mathcal{E}_1-\mathcal{E}_2)}{R_1\mathcal{E}_2}
}
\]

ถ้า $\mathcal{E}_1<\mathcal{E}_2$ ค่า $i_1$ เป็นลบ แสดงว่ากระแสจริงสวนทางกับลูกศรที่กำหนด
\end{examplebox}

% =========================================================
\subsection{ปี 2565 ข้อ 2: กระแสในวงจรหลายแขนงที่มีแหล่งกำเนิดต่างกัน}

\vspace{1em}

\IfFileExists{img/posn65q2.png}{
\begin{center}
    \includegraphics[width=1\linewidth]{img/posn65q2.png}
\end{center}
}{
\begin{tcolorbox}[title=ตำแหน่งภาพที่แนะนำ]
ให้ตัดภาพข้อสอบปี 2565 ข้อ 2 แล้วบันทึกเป็น
\[
\texttt{img/posn65q2.png}
\]
\end{tcolorbox}
}

\begin{examplebox}{เฉลยละเอียด}
วงจรมีสี่แขนงขนานระหว่างโหนดซ้ายและขวา แต่ละแขนงมีตัวต้านทาน $R$

ให้

\[
V\equiv V_{\text{left}}-V_{\text{right}}
\]

และกำหนดกระแสจากซ้ายไปขวาเป็นบวก

สามแขนงบนมีแรงเคลื่อนไฟฟ้า $\mathcal{E}$ จึงมีกระแสแต่ละแขนง

\[
I_{\mathcal{E}}=\frac{V-\mathcal{E}}{R}
\]

แขนงล่างมีแรงเคลื่อนไฟฟ้า $2\mathcal{E}$

\[
I_{2\mathcal{E}}=\frac{V-2\mathcal{E}}{R}
\]

เพราะไม่มีแขนงภายนอกต่อเข้ากับโหนด ใช้กฎจุดต่อ

\[
3\left(\frac{V-\mathcal{E}}{R}\right)
+\frac{V-2\mathcal{E}}{R}=0
\]

\[
4V-5\mathcal{E}=0
\]

\[
V=\frac{5\mathcal{E}}{4}
\]

กระแสผ่านจุด $A$ อยู่ในหนึ่งในสามแขนงบน

\[
I_A=\frac{V-\mathcal{E}}{R}
\]

\[
I_A=\frac{5\mathcal{E}/4-\mathcal{E}}{R}
\]

\[
\boxed{I_A=\frac{\mathcal{E}}{4R}}
\]

ตอบ ง
\end{examplebox}

% ---------------------------------------------------------
\subsection{ปี 2565 ข้อ 6: อัตราการทำงานของแหล่งกำเนิดขณะชาร์จตัวเก็บประจุ}

\vspace{1em}

\IfFileExists{img/posn65q6.png}{
\begin{center}
    \includegraphics[width=1\linewidth]{img/posn65q6.png}
\end{center}
}{
\begin{tcolorbox}[title=ตำแหน่งภาพที่แนะนำ]
ให้ตัดภาพข้อสอบปี 2565 ข้อ 6 แล้วบันทึกเป็น
\[
\texttt{img/posn65q6.png}
\]
\end{tcolorbox}
}

\begin{examplebox}{เฉลยละเอียด}
ใช้กฎลูปในขณะตัวเก็บประจุมีประจุ $q$ และมีกระแส $I$

\[
\mathcal{E}-IR-\frac{q}{C}=0
\]

ดังนั้น

\[
\mathcal{E}=IR+\frac{q}{C}
\]

กำลังที่แหล่งกำเนิดจ่ายคือ

\[
P_{\text{source}}=\mathcal{E}I
\]

คูณสมการลูปด้วย $I$

\[
\mathcal{E}I=RI^2+\frac{qI}{C}
\]

จึงได้

\[
\boxed{P_{\text{source}}=RI^2+\frac{qI}{C}}
\]

พจน์แรกเป็นกำลังที่เปลี่ยนเป็นความร้อนในตัวต้านทาน ส่วนพจน์ที่สองเป็นอัตราการเพิ่มพลังงานในสนามไฟฟ้าของตัวเก็บประจุ

ตอบ ค
\end{examplebox}

% =========================================================
\subsection{ปี 2566 ข้อ 13: รูปร่างกราฟประจุในวงจรตัวเก็บประจุสองตัว}

\vspace{1em}

\IfFileExists{img/posn66q13.png}{
\begin{center}
    \includegraphics[width=1\linewidth]{img/posn66q13.png}
\end{center}
}{
\begin{tcolorbox}[title=ตำแหน่งภาพที่แนะนำ]
ให้ตัดภาพข้อสอบปี 2566 ข้อ 13 แล้วบันทึกเป็น
\[
\texttt{img/posn66q13.png}
\]
\end{tcolorbox}
}

\begin{examplebox}{เฉลยละเอียด}
ตัวเก็บประจุทั้งสองเริ่มไม่มีประจุ ดังนั้นตัวเก็บประจุด้านขวาที่มีเครื่องหมาย $*$ มี

\[
q_*(0)=0
\]

ทันทีหลังปิดสวิตช์ แหล่งกำเนิดทำให้เกิดกระแสชั่วคราว ตัวเก็บประจุด้านขวาจึงเริ่มสะสมประจุในทิศหนึ่ง ทำให้ $q_*$ เพิ่มจากศูนย์

เมื่อเวลานานมาก ตัวเก็บประจุตัวบนที่ต่ออนุกรมกับแหล่งกำเนิดทำให้กระแส DC ผ่านแขนงนั้นเป็นศูนย์ ส่วนตัวเก็บประจุด้านขวาต่อขนานกับตัวต้านทาน $R$ จึงคายประจุผ่านตัวต้านทานจนแรงดันคร่อมเป็นศูนย์

\[
q_*(\infty)=0
\]

วงจรมีเพียงตัวต้านทานและตัวเก็บประจุ ไม่มีตัวเหนี่ยวนำ จึงไม่มีการสั่นกลับไปกลับมา ประจุต้องเพิ่มขึ้นถึงค่าสูงสุดหนึ่งครั้ง แล้วลดกลับเข้าหาศูนย์

ดังนั้นกราฟที่ถูกต้องคือ

\[
\boxed{\text{กราฟ 2}}
\]
\end{examplebox}

% ---------------------------------------------------------
\subsection{ปี 2566 ข้อ 19: วงจรสมมาตรที่มีตัวต้านทานเท่ากันทุกตัว}

\vspace{1em}

\IfFileExists{img/posn66q19.png}{
\begin{center}
    \includegraphics[width=1\linewidth]{img/posn66q19.png}
\end{center}
}{
\begin{tcolorbox}[title=ตำแหน่งภาพที่แนะนำ]
ให้ตัดภาพข้อสอบปี 2566 ข้อ 19 แล้วบันทึกเป็น
\[
\texttt{img/posn66q19.png}
\]
\end{tcolorbox}
}

\begin{examplebox}{เฉลยละเอียด}
วงจรประกอบด้วยสี่โหนด และมีตัวต้านทาน $R$ เชื่อมทุกคู่ของโหนด จึงมีสมมาตรเมื่อมองจากปลาย $A$ และ $B$

โหนดภายในสองโหนดต้องมีศักย์เท่ากัน จึงไม่มีกระแสผ่านตัวต้านทานที่เชื่อมระหว่างโหนดภายในทั้งสอง

เมื่อนำโหนดภายในที่ศักย์เท่ากันมารวมเป็นโหนดเดียว

จาก $A$ ไปโหนดกลางมีตัวต้านทาน $R$ สองตัวขนานกัน

\[
R_{A\to M}=R\parallel R=\frac{R}{2}
\]

จากโหนดกลางไป $B$ ก็มี

\[
R_{M\to B}=\frac{R}{2}
\]

เส้นทางผ่านโหนดกลางจึงมี

\[
R_{\text{path}}=\frac{R}{2}+\frac{R}{2}=R
\]

และยังมีตัวต้านทาน $R$ เชื่อมตรงจาก $A$ ไป $B$ อีกหนึ่งแขนง

ดังนั้น

\[
R_{AB}=R\parallel R
\]

\[
\boxed{R_{AB}=\frac{R}{2}}
\]

ตอบข้อ 2
\end{examplebox}

% =========================================================
\subsection{ปี 2567 ข้อ 7: ตัวเก็บประจุหลายแผ่นต่อสลับขั้ว}

\vspace{1em}

\IfFileExists{img/posn67q7.png}{
\begin{center}
    \includegraphics[width=1\linewidth]{img/posn67q7.png}
\end{center}
}{
\begin{tcolorbox}[title=ตำแหน่งภาพที่แนะนำ]
ให้ตัดภาพข้อสอบปี 2567 ข้อ 7 แล้วบันทึกเป็น
\[
\texttt{img/posn67q7.png}
\]
\end{tcolorbox}
}

\begin{examplebox}{เฉลยละเอียด}
แผ่นโลหะสี่แผ่นวางสลับขั้วกัน ทำให้ช่องว่างระหว่างแผ่นที่อยู่ติดกันทั้งสามช่องมีความต่างศักย์เท่ากับ $V$

แต่ละช่องทำตัวเป็นตัวเก็บประจุแผ่นขนาน

\[
C_0=\varepsilon_0\frac{A}{d}
\]

ตัวเก็บประจุทั้งสามช่องต่อขนาน เพราะต่ออยู่ระหว่างโหนดสองโหนดเดียวกัน

\[
C_{\text{eq}}=C_0+C_0+C_0
\]

\[
C_{\text{eq}}=3\varepsilon_0\frac{A}{d}
\]

ขนาดประจุสุทธิที่แหล่งกำเนิดส่งให้ระบบคือ

\[
Q=C_{\text{eq}}V
\]

ดังนั้น

\[
\boxed{Q=\frac{3\varepsilon_0AV}{d}}
\]
\end{examplebox}

% ---------------------------------------------------------
\subsection{ปี 2567 ข้อ 8: ประจุในตัวเก็บประจุที่ขนานกับตัวแบ่งแรงดัน}

\vspace{1em}

\IfFileExists{img/posn67q8.png}{
\begin{center}
    \includegraphics[width=1\linewidth]{img/posn67q8.png}
\end{center}
}{
\begin{tcolorbox}[title=ตำแหน่งภาพที่แนะนำ]
ให้ตัดภาพข้อสอบปี 2567 ข้อ 8 แล้วบันทึกเป็น
\[
\texttt{img/posn67q8.png}
\]
\end{tcolorbox}
}

\begin{examplebox}{เฉลยละเอียด}
เมื่อกระแสในวงจรคงที่แล้ว ตัวเก็บประจุไม่มีกระแสผ่าน แต่แรงดันคร่อมตัวเก็บประจุเท่ากับแรงดันคร่อมตัวต้านทานที่ต่อขนานกัน

ดังนั้น

\[
V_{C_1}=IR_1
\]

และ

\[
V_{C_2}=IR_2
\]

ประจุคือ

\[
q_1=C_1V_{C_1}=C_1IR_1
\]

\[
q_2=C_2V_{C_2}=C_2IR_2
\]

จึงได้

\[
\frac{q_1}{q_2}
=\frac{R_1C_1}{R_2C_2}
\]

โจทย์กำหนดว่า

\[
R_1C_1=R_2C_2
\]

ดังนั้น

\[
\boxed{\frac{q_1}{q_2}=1}
\]
\end{examplebox}

% ---------------------------------------------------------
\subsection{ปี 2567 ข้อ 16: ความจุสมมูลและศักย์ในวงจรบันไดตัวเก็บประจุ}

\vspace{1em}

\IfFileExists{img/posn67q16.png}{
\begin{center}
    \includegraphics[width=1\linewidth]{img/posn67q16.png}
\end{center}
}{
\begin{tcolorbox}[title=ตำแหน่งภาพที่แนะนำ]
ให้ตัดภาพข้อสอบปี 2567 ข้อ 16 แล้วบันทึกเป็น
\[
\texttt{img/posn67q16.png}
\]
\end{tcolorbox}
}

\begin{examplebox}{เฉลยละเอียด}
\textbf{ก. หาค่า $C$ ที่ทำให้ $C_{AB}=C$}

ตัวเก็บประจุ $C_2$ และ $C$ ต่อขนาน

\[
C_p=C_2+C
\]

ต่ออนุกรมกับ $C_1$

\[
C_{AB}=\frac{C_1(C_2+C)}{C_1+C_2+C}
\]

กำหนดให้ $C_{AB}=C$

\[
C(C_1+C_2+C)=C_1(C_2+C)
\]

จัดรูป

\[
C^2+C_2C-C_1C_2=0
\]

เลือกรากบวก

\[
\boxed{
C=\frac{\sqrt{C_2^2+4C_1C_2}-C_2}{2}
}
\]

\textbf{ข. หาศักย์ไฟฟ้าที่จุด 5}

ลดวงจรจากด้านขวา

ที่จุด 5 มีตัวเก็บประจุ $4\,\si{\micro\farad}$ ต่อกราวด์ จึงมีความจุสมมูลด้านขวาสุด

\[
C_{5}=4\,\si{\micro\farad}
\]

มองจากจุด 4 ตัวเก็บประจุอนุกรม $4\,\si{\micro\farad}$ กับ $C_5=4\,\si{\micro\farad}$ ให้

\[
C_s=\frac{(4)(4)}{4+4}=2\,\si{\micro\farad}
\]

ขนานกับตัวเก็บประจุแนวดิ่ง $2\,\si{\micro\farad}$ ที่จุด 4

\[
C_{4}=2+2=4\,\si{\micro\farad}
\]

โครงสร้างเดียวกันเกิดซ้ำที่จุด 3, 2 และ 1 จึงมีความจุสมมูลจากแต่ละจุดลงกราวด์เท่ากับ $4\,\si{\micro\farad}$

ระหว่างแหล่งจ่ายกับจุด 1 จึงเป็นตัวเก็บประจุ $4\,\si{\micro\farad}$ สองส่วนอนุกรมเท่ากัน ทำให้

\[
V_1=\frac{\mathcal{E}}{2}
\]

ระหว่างจุด 1 กับจุด 2 ก็เป็นการแบ่งแรงดันระหว่างความจุเท่ากันอีกครั้ง

\[
V_2=\frac{V_1}{2}=\frac{\mathcal{E}}{4}
\]

ดำเนินต่อไป

\[
V_n=\frac{\mathcal{E}}{2^n}
\]

ดังนั้น

\[
\boxed{V_5=\frac{\mathcal{E}}{32}}
\]
\end{examplebox}

% =========================================================
\subsection{ปี 2568 ข้อ 7: แบ่งลวดแล้วต่อขนานให้ได้ความต้านทานสูงสุด}

\vspace{1em}

\IfFileExists{img/posn68q7.png}{
\begin{center}
    \includegraphics[width=1\linewidth]{img/posn68q7.png}
\end{center}
}{
\begin{tcolorbox}[title=ตำแหน่งภาพที่แนะนำ]
ให้ตัดภาพข้อสอบปี 2568 ข้อ 7 แล้วบันทึกเป็น
\[
\texttt{img/posn68q7.png}
\]
\end{tcolorbox}
}

\begin{examplebox}{เฉลยละเอียด}
ลวดเดิมมีความต้านทานรวม $R$ เมื่อแบ่งเป็นสองส่วน

\[
R_1+R_2=R
\]

นำสองส่วนมาต่อขนาน

\[
R_{\parallel}=\frac{R_1R_2}{R_1+R_2}
\]

ใช้ $R_1+R_2=R$

\[
R_{\parallel}=\frac{R_1R_2}{R}
\]

สำหรับผลบวกคงตัว ผลคูณมากที่สุดเมื่อสองจำนวนเท่ากัน

\[
R_1=R_2=\frac{R}{2}
\]

ดังนั้น

\[
R_{\parallel,\max}
=\frac{(R/2)(R/2)}{R}
\]

\[
\boxed{R_{\parallel,\max}=\frac{R}{4}}
\]
\end{examplebox}

% ---------------------------------------------------------
\subsection{ปี 2568 ข้อ 8: ตัวต้านทานที่ทำให้ความต้านทานรวมเท่ากับตัวมันเอง}

\vspace{1em}

\IfFileExists{img/posn68q8.png}{
\begin{center}
    \includegraphics[width=1\linewidth]{img/posn68q8.png}
\end{center}
}{
\begin{tcolorbox}[title=ตำแหน่งภาพที่แนะนำ]
ให้ตัดภาพข้อสอบปี 2568 ข้อ 8 แล้วบันทึกเป็น
\[
\texttt{img/posn68q8.png}
\]
\end{tcolorbox}
}

\begin{examplebox}{เฉลยละเอียด}
วงจรมี $R_1$ ต่ออนุกรมกับวงขนานของ $R_2$ และตัวต้านทานที่ต้องหา $R$

\[
R_{AB}=R_1+\frac{R_2R}{R_2+R}
\]

โจทย์กำหนดว่า

\[
R_{AB}=R
\]

จึงได้

\[
R=R_1+\frac{R_2R}{R_2+R}
\]

คูณด้วย $R_2+R$

\[
R(R_2+R)=R_1(R_2+R)+R_2R
\]

ตัดพจน์ $R_2R$

\[
R^2=R_1R_2+R_1R
\]

\[
R^2-R_1R-R_1R_2=0
\]

ใช้สูตรสมการกำลังสองและเลือกรากบวก

\[
\boxed{
R=\frac{R_1+\sqrt{R_1^2+4R_1R_2}}{2}
}
\]
\end{examplebox}

% ---------------------------------------------------------
\subsection{ปี 2568 ข้อ 16: การอนุรักษ์กำลังระหว่างชาร์จตัวเก็บประจุ}

\vspace{1em}

\IfFileExists{img/posn68q16.png}{
\begin{center}
    \includegraphics[width=1\linewidth]{img/posn68q16.png}
\end{center}
}{
\begin{tcolorbox}[title=ตำแหน่งภาพที่แนะนำ]
ให้ตัดภาพข้อสอบปี 2568 ข้อ 16 แล้วบันทึกเป็น
\[
\texttt{img/posn68q16.png}
\]
\end{tcolorbox}
}

\begin{examplebox}{เฉลยละเอียด}
ใช้กฎลูป

\[
\mathcal{E}=iR+\frac{q}{C}
\]

คูณทั้งสมการด้วย $i$

\[
\mathcal{E}i=i^2R+\frac{q}{C}i
\]

ดังนั้นปริมาณที่โจทย์ถามคือ

\[
\boxed{i^2R+\frac{q}{C}i=\mathcal{E}i}
\]

ด้านซ้ายคือผลรวมของกำลังที่เปลี่ยนเป็นความร้อนในตัวต้านทานกับอัตราการเพิ่มพลังงานในตัวเก็บประจุ ส่วนด้านขวาคือกำลังที่แหล่งกำเนิดจ่าย
\end{examplebox}

% ---------------------------------------------------------
\subsection{ปี 2568 ข้อ 17: ตัวเก็บประจุอนุกรมจำนวนมาก}

\vspace{1em}

\IfFileExists{img/posn68q17.png}{
\begin{center}
    \includegraphics[width=1\linewidth]{img/posn68q17.png}
\end{center}
}{
\begin{tcolorbox}[title=ตำแหน่งภาพที่แนะนำ]
ให้ตัดภาพข้อสอบปี 2568 ข้อ 17 แล้วบันทึกเป็น
\[
\texttt{img/posn68q17.png}
\]
\end{tcolorbox}
}

\begin{examplebox}{เฉลยละเอียด}
ตัวเก็บประจุทั้งหมดต่ออนุกรม จึงมีขนาดประจุเท่ากันทุกตัว

โจทย์กำหนด

\[
C_i=i^2C_1
\]

ความจุสมมูลของชุดอนุกรมคือ

\[
\frac{1}{C_{\text{eq}}}
=\sum_{i=1}^{n}\frac{1}{C_i}
\]

แทน $C_i=i^2C_1$

\[
\frac{1}{C_{\text{eq}}}
=\frac{1}{C_1}\sum_{i=1}^{n}\frac{1}{i^2}
\]

เมื่อ $n$ มีค่ามาก โจทย์ให้

\[
\sum_{i=1}^{n}\frac{1}{i^2}
\approx\frac{\pi^2}{6}
\]

ดังนั้น

\[
\frac{1}{C_{\text{eq}}}
\approx\frac{\pi^2}{6C_1}
\]

\[
C_{\text{eq}}\approx\frac{6C_1}{\pi^2}
\]

ประจุในวงจรอนุกรมคือ

\[
q=C_{\text{eq}}\mathcal{E}
\]

จึงมี

\[
\boxed{q_i\approx\frac{6C_1\mathcal{E}}{\pi^2}}
\]

สำหรับตัวเก็บประจุทุกตัว รวมถึง $C_i$
\end{examplebox}

% =========================================================
% End of chapter
% =========================================================
